\documentclass[11pt]{article}

\usepackage[a4paper,margin=1in]{geometry}
\usepackage{amsmath,amssymb,amsthm,mathtools}
\usepackage{enumitem}
\usepackage[hidelinks]{hyperref}

\setlength{\parskip}{0.3em}

\newtheorem{theorem}{Theorem}[section]
\newtheorem{proposition}[theorem]{Proposition}
\newtheorem{lemma}[theorem]{Lemma}
\newtheorem{corollary}[theorem]{Corollary}
\theoremstyle{definition}
\newtheorem{definition}[theorem]{Definition}
\newtheorem{remark}[theorem]{Remark}

\newcommand{\R}{\mathbb{R}}
\newcommand{\Om}{\Omega}
\newcommand{\HH}{\mathcal{H}}
\newcommand{\eps}{\varepsilon}
\newcommand{\dT}{\delta_T}
\newcommand{\diso}{\delta_{\mathrm{iso}}}
\newcommand{\Etil}{\widetilde E}
\newcommand{\abs}[1]{\left|#1\right|}

\title{The inset of a near-extremal Saint--Venant domain\\
       is a regular set}
\author{}
\date{}

\begin{document}
\maketitle

\begin{abstract}
Let $\Om\subset\R^2$ be a bounded open set of area $\pi$ whose torsional
rigidity is close to that of the disk, i.e.\ with small Saint--Venant
deficit $\dT(\Om)=T(B_1)-T(\Om)$. The boundary $\partial\Om$ may be
arbitrarily irregular. We prove that a suitably chosen super-level set of
the torsion function --- the \emph{inset} --- is nonetheless a regular set,
and we quantify in what sense. Precisely (Theorem~\ref{thm:main}), there is
a regular value $\hat v$ of the torsion function $u$ such that the inset
$\Etil=\{u>\hat v\}$ has: \emph{(i)} an analytic boundary (a finite union of
analytic Jordan curves) lying at positive distance from $\partial\Om$, with
$|\nabla u|>0$ on it; \emph{(ii)} small isoperimetric and Serrin deficits,
both $O(\dT)$; \emph{(iii)} closeness to a disk in measure,
$|\Etil\,\triangle\,B|\lesssim\sqrt\dT$; \emph{(iv)} no interior holes; and
\emph{(v)} containment in a thin annulus between two concentric circles,
$B_{\hat r-\eta}(z)\subset\Etil\subset B_{\hat r+\eta}(z)$. For the width we
give two routes: applying the Magnanini--Poggesi stability theorem to the
\emph{smooth} inset yields $\eta\lesssim\dT^{\tau/2}$ but with a constant
depending on the inset's $C^2$ geometry; replacing the interior/exterior
ball conditions of that proof by the \emph{measure-theoretic} crossing count
of the radial function on $\partial\Etil$ --- the coarea identity
$\int_{\partial\Etil}|\nu_\theta|=\int_0^\infty N(\rho)\,d\rho$, which sums
all inward fjords and outward tentacles at once --- yields instead
$\eta\lesssim\dT^{1/4}$ with an \emph{absolute} constant, free of any $C^2$
dependence and free of the multi-fjord covering bookkeeping, the whole
dependence being on the isoperimetric deficit $D_I$. Every lemma is proved
in full; standard facts are cited verbatim. The qualitative principle is
that the roughness of $\partial\Om$ never reaches the inset: every feature
of $\Etil$ is controlled by the deficits alone.
\end{abstract}

\tableofcontents

\section{Introduction}\label{sec:intro}

\subsection*{Setup and motivation}
For a bounded open set $\Om\subset\R^2$, the \emph{torsion function} $u$ is
the solution of $-\Delta u=1$ in $\Om$ with $u=0$ on $\partial\Om$, and the
\emph{torsional rigidity} is $T(\Om)=\int_\Om u$. The Saint--Venant
inequality states $T(\Om)\le T(B)$ among sets $B$ of equal area, with
equality only for disks. Normalising $|\Om|=\pi$ (so the optimal disk is the
unit disk $B_1$, $T(B_1)=\pi/8$), the \emph{Saint--Venant deficit} is
\[
   \dT(\Om):=T(B_1)-T(\Om)\ \ge\ 0 .
\]
A central problem is the sharp \emph{quantitative} Saint--Venant
inequality, which would bound the Fraenkel asymmetry of $\Om$ by $\dT$. A
standard difficulty is that $\partial\Om$ carries no a priori regularity:
the only hypothesis is that $\dT$ is small. The contribution of this note is
to show, with complete proofs, that one can extract from $\Om$ a
\emph{regular} subdomain --- a super-level set of $u$ --- that captures the
geometry while being analytic, nearly circular, and trapped (on the outside)
in a thin annulus. This ``inset'' is the object on which the perturbative
machinery of the quantitative theory can be run.

The key principle is that \emph{the interior is always regular}: $u$ solves
an analytic elliptic equation, so it is real-analytic inside $\Om$ no matter
how rough $\partial\Om$ is, and its level sets a definite distance inside are
analytic curves (Section~\ref{sec:smooth}). The roughness of $\partial\Om$
does not propagate inward. What requires work is the \emph{quantitative}
geometry: that the inset is nearly a disk, and how nearly.

\subsection*{Standing assumptions}
Throughout, $\Om\subset\R^2$ is a bounded open set, regular for the
Dirichlet problem (so that the torsion function satisfies $u\in
H^1_0(\Om)\cap C(\overline\Om)$ with $u=0$ on $\partial\Om$; this holds for
instance whenever $\Om$ satisfies an exterior corkscrew or cone condition,
and is no loss for the deficit since $\dT$ is continuous under invasion of
$\Om$ by regular subdomains). We normalise $|\Om|=\pi$. We write $\dT$ for
$\dT(\Om)$ and assume $\dT\le\delta_\star$, where $\delta_\star>0$ is an
absolute constant fixed in the course of the proofs. Constants denoted
$c,C,C_0,\dots$ are absolute unless a dependence is displayed; the only
non-absolute dependence that appears anywhere is in the annulus
\emph{width} of Section~\ref{sec:trapping}, whose constant depends on the
$C^2$ geometry of the (smooth) inset --- see Remark~\ref{rem:nonuniform}.
Every other constant in the paper is absolute.

\subsection*{Main result}
The following theorem collects the conclusions proved in
Sections~\ref{sec:smooth}--\ref{sec:trapping}. Recall the \emph{isoperimetric
deficit} $\diso(E)=P(E)/(2\sqrt{\pi|E|})-1$ and the \emph{Fraenkel
asymmetry} $\mathcal A(E)=\min_x|E\triangle B(x)|/|E|$ (disk $B(x)$ of area
$|E|$).

\begin{theorem}[The inset is a regular set]\label{thm:main}
There is an absolute constant $\delta_\star>0$ with the following property.
Let $\Om\subset\R^2$ be bounded open, regular for the Dirichlet problem,
$|\Om|=\pi$, with torsion function $u$ and $\dT(\Om)\le\delta_\star$. Then
there is a regular value $\hat v\in(0,\tfrac14)$ of $u$ such that the
\emph{inset}
\[
   \Etil:=\{x\in\Om:u(x)>\hat v\},\qquad
   \hat\mu:=|\Etil|\in[\tfrac\pi4,\pi),\qquad
   \hat r:=\sqrt{\hat\mu/\pi}\in[\tfrac12,1),
\]
satisfies all of the following.
\begin{enumerate}[label=\textup{(\roman*)},leftmargin=2.4em]
\item \textup{(Analytic boundary, interior.)} $\partial\Etil=\{u=\hat v\}$
is a finite disjoint union of real-analytic Jordan curves, with
$|\nabla u|>0$ everywhere on it, and
$\operatorname{dist}(\overline\Etil,\partial\Om)>0$. Consequently $\Etil$ is
a bounded open set of finite perimeter, $P(\Etil)=\HH^1(\{u=\hat v\})$, with
outward unit normal $-\nabla u/|\nabla u|$.
\item \textup{(Controlled deficits.)} $\diso(\Etil)\le C\,\dT$ and the
weighted Serrin deficit obeys
$\displaystyle\int_{\partial\Etil}\frac{(|\nabla u|-c)^2}{|\nabla u|}\,
d\HH^1\le C\,\dT$, where $c=|\Etil|/P(\Etil)$.
\item \textup{(Near-circular in measure.)} There is $z\in\R^2$ with
$\abs{\Etil\,\triangle\,B_{\hat r}(z)}\le C\,\sqrt\dT$; equivalently
$\mathcal A(\Etil)\le C\sqrt\dT$.
\item \textup{(No interior holes.)} Every connected component of
$\Om\setminus\overline\Etil$ has closure meeting $\partial\Om$; in
particular $\Etil$ has no holes.
\item \textup{(Squeezed between two concentric circles.)} The inset (its
main component) is trapped in a thin annulus
$B_{\hat r-\eta}(z)\subset\Etil\subset B_{\hat r+\eta}(z)$, in either of two
forms:
\begin{itemize}[leftmargin=1.4em]
\item[\textup{(v$_{\mathrm{a}}$)}] \textup{(via Magnanini--Poggesi)}
$\eta\le C_{\mathrm{MP}}\,\dT^{\tau/2}$, $\tau\in(0,1]$ absolute, but
$C_{\mathrm{MP}}$ depending on the $C^2$ geometry of the smooth inset
\textup{(}Theorem~\ref{thm:annulus}\textup{)};
\item[\textup{(v$_{\mathrm{b}}$)}] \textup{(regularity-free)} $\eta\le
C\,\dT^{1/4}$ with $C$ \emph{absolute}, depending only on the deficits;
proved via the angular-perimeter coarea identity
$P_\theta=\int_0^\infty N(\rho)\,d\rho$, which charges all inward fjords and
outward tentacles to $D_I$ at once \textup{(}no $C^2$ data, no covering;
Theorem~\ref{thm:perim-annulus}\textup{)}. The only standing assumption is
connectedness of the main component, which carries area $\hat\mu-O(\dT)$.
\end{itemize}
\end{enumerate}
The constants $C$ in \textup{(i)--(iv)} are \emph{absolute}. Clause
\textup{(v$_{\mathrm a}$)} pays a non-absolute $C^2$ constant;
\textup{(v$_{\mathrm b}$)} removes it with an absolute constant, the entire
dependence then being on the isoperimetric deficit $D_I$
(Remarks~\ref{rem:nonuniform}, \ref{rem:covering-gone}, \ref{rem:loop}).
\end{theorem}

\subsection*{Strategy and organisation}
The proof has three layers. \emph{Interior regularity}
(Sections~\ref{sec:prelim}--\ref{sec:smooth}): $u$ is analytic inside, so
a.e.\ level set is an analytic curve at positive distance from
$\partial\Om$; this is qualitative and uses no smallness. \emph{Energy /
$L^1$ control} (Sections~\ref{sec:identity}--\ref{sec:controlled}): the
level-set deficit identity
$\int_0^m(D_H+D_I)\,dv=4\pi\dT$ (Theorem~\ref{thm:identity}) writes the
deficit as the total of two non-negative level defects; combined with
Talenti's volume comparison (Section~\ref{sec:talenti}) it yields, by
averaging, a single good inset with $O(\dT)$ isoperimetric and Serrin
defects and definite area, hence $\sqrt\dT$-closeness to a disk by the
sharp quantitative isoperimetric inequality. \emph{Geometric control}
(Sections~\ref{sec:topology}--\ref{sec:trapping}): superharmonicity removes
interior holes, and --- crucially using that the inset is \emph{smooth} ---
the Magnanini--Poggesi stability theorem for Serrin's overdetermined problem
traps the inset between two concentric circles. This last step is where the
inset's smoothness pays off: the regularity-dependent stability machinery,
inapplicable to the rough $\Om$, applies to the analytic $\Etil$.

Section~\ref{sec:prelim} collects the cited background; the reader familiar
with finite-perimeter theory, the FMP inequality, interior elliptic
regularity, and Sard's and Talenti's theorems may skip to
Section~\ref{sec:smooth}.

\bigskip
\section{Preliminaries}\label{sec:prelim}

Throughout this paper $\Om\subset\R^2$ denotes a bounded open set, $|\cdot|$
denotes Lebesgue measure on $\R^2$ (or $\R^n$ where stated), $\HH^1$ denotes
$1$-dimensional Hausdorff measure, and $B_r(x)\subset\R^2$ denotes the open
disk of radius $r$ centred at $x$, with the shorthand $B_r=B_r(0)$. We write
$\omega_n=|B_1|$ for the volume of the unit ball in $\R^n$, so that
$\omega_2=\pi$. We collect here the definitions and results that the rest of
the paper relies on. Standard results are stated precisely and cited verbatim
to published sources; the few statements that are not direct citations are
proved.

\subsection{Sets of finite perimeter}\label{ss:perimeter}

We follow the conventions of Maggi's monograph.

\begin{definition}[Perimeter]\label{def:perimeter}
Let $E\subset\R^2$ be a Lebesgue measurable set and let $A\subset\R^2$ be
open. The \emph{perimeter of $E$ in $A$} is
\[
  P(E;A):=\sup\Big\{\int_E \operatorname{div}\varphi\,dx
    \;:\;\varphi\in C^1_c(A;\R^2),\ \|\varphi\|_{L^\infty(A)}\le 1\Big\}.
\]
We write $P(E):=P(E;\R^2)$ for the (total) perimeter. The set $E$ is said to
be of \emph{finite perimeter} (in $A$) if $P(E;A)<\infty$.
\end{definition}

If $E$ has finite perimeter in $A$, then $\mu_E:=-D\mathbf 1_E$, the
distributional gradient of the indicator $\mathbf 1_E$, is an $\R^2$-valued
Radon measure on $A$ with total variation $|\mu_E|(A)=P(E;A)$; see
\cite[Proposition~12.1 and Definition~12.1]{Maggi2012}.

\begin{definition}[Reduced boundary]\label{def:reduced-boundary}
Let $E\subset\R^2$ have locally finite perimeter and let $\mu_E=-D\mathbf 1_E$
be the associated Gauss--Green measure. The \emph{reduced boundary}
$\partial^*E$ is the set of points $x\in\operatorname{supp}\mu_E$ for which the
limit
\[
  \nu_E(x):=\lim_{r\to0^+}\frac{\mu_E(B_r(x))}{|\mu_E|(B_r(x))}
\]
exists and satisfies $|\nu_E(x)|=1$. The vector $\nu_E(x)$ is the
\emph{measure-theoretic outer unit normal} to $E$ at $x$.
\end{definition}

\begin{theorem}[De Giorgi structure theorem]\label{thm:de-giorgi}
Let $E\subset\R^2$ be a set of locally finite perimeter. Then $\partial^*E$ is
countably $\HH^1$-rectifiable, the Gauss--Green measure satisfies
$\mu_E=\nu_E\,\HH^1\llcorner\partial^*E$, and consequently
\[
  P(E;A)=\HH^1(\partial^*E\cap A)\qquad\text{for every Borel set }A\subset\R^2 .
\]
\end{theorem}

\begin{proof}
This is De Giorgi's structure theorem; see
\cite[Theorem~15.9]{Maggi2012} together with
\cite[Corollary~16.1]{Maggi2012}, or
\cite[Section~5.7.3, Theorem~5.15]{EvansGariepy2015}.
\end{proof}

For $x\in\R^2$ and $t\in[0,1]$ we say that $E$ has \emph{density $t$ at $x$} if
\[
  \Theta(E,x):=\lim_{r\to0^+}\frac{|E\cap B_r(x)|}{|B_r(x)|}=t .
\]
We write $E^{(t)}:=\{x\in\R^2:\Theta(E,x)=t\}$. The sets $E^{(1)}$ and
$E^{(0)}$ are the \emph{measure-theoretic interior} and \emph{exterior} of
$E$, respectively.

\begin{definition}[Essential boundary]\label{def:essential-boundary}
The \emph{essential boundary} of $E$ is
\[
  \partial^e E:=\R^2\setminus\big(E^{(0)}\cup E^{(1)}\big),
\]
the set of points at which $E$ has neither density $0$ nor density $1$.
\end{definition}

\begin{theorem}[Federer; density $1/2$ on the reduced boundary]\label{thm:federer}
Let $E\subset\R^2$ be a set of locally finite perimeter. Then
\[
  \partial^*E\subseteq E^{(1/2)}\subseteq\partial^e E,
  \qquad
  \HH^1\big(\partial^e E\setminus\partial^*E\big)=0 .
\]
In particular $E$ has density $\tfrac12$ at $\HH^1$-almost every point of
$\partial^*E$, and
$P(E;A)=\HH^1(\partial^*E\cap A)=\HH^1(\partial^e E\cap A)$ for every Borel
set $A$.
\end{theorem}

\begin{proof}
The inclusion $\partial^*E\subseteq E^{(1/2)}$ is part of the blow-up analysis
of the reduced boundary, \cite[Theorem~15.5]{Maggi2012}. Federer's theorem
$\HH^1(\partial^e E\setminus\partial^*E)=0$ is
\cite[Theorem~16.2]{Maggi2012}; see also
\cite[Section~5.8, Theorem~5.16 (Federer's theorem)]{EvansGariepy2015}.
\end{proof}

\subsection{The coarea formula}\label{ss:coarea}

\begin{theorem}[Coarea formula]\label{thm:coarea}
Let $u\in W^{1,1}(\Om)$ (in particular, any Lipschitz function $u$ on $\Om$).
Then for every nonnegative Borel function $g:\Om\to[0,\infty]$,
\[
  \int_\Om g(x)\,|\nabla u(x)|\,dx
  =\int_{-\infty}^{\infty}\Big(\int_{\{u=t\}}g\,d\HH^1\Big)\,dt,
\]
where $\{u=t\}:=\{x\in\Om:u(x)=t\}$ and $u$ is identified with a precise
representative. Moreover, for Lebesgue-almost every $t\in\R$ the level set
$\{u=t\}$ is $\HH^1$-rectifiable, $\HH^1(\{u=t\})<\infty$, and the
superlevel set $\{u>t\}$ has finite perimeter in $\Om$ with
\[
  \partial^*\{u>t\}\cap\Om\subseteq\{u=t\},
  \qquad
  P(\{u>t\};\Om)=\HH^1(\{u=t\})\quad\text{for a.e.\ }t .
\]
\end{theorem}

\begin{proof}
The coarea formula for Sobolev (and Lipschitz) functions is
\cite[Theorem~13.1 and Remark~13.2]{Maggi2012}; for the Lipschitz case see
also \cite[Section~3.4.3, Theorem~3.13]{EvansGariepy2015}. The
rectifiability of almost every level set and the identity
$P(\{u>t\};\Om)=\HH^1(\{u=t\})$ for a.e.\ $t$, together with
$\partial^*\{u>t\}\cap\Om\subseteq\{u=t\}$, follow from the coarea formula for
$BV$ functions \cite[Theorem~3.40]{AFP2000}: indeed
$\int_\R P(\{u>t\};\Om)\,dt=\int_\Om|\nabla u|\,dx<\infty$, so $\{u>t\}$ has
finite perimeter for a.e.\ $t$, and the inclusion of the reduced boundary in
the level set is \cite[Theorem~16.2]{Maggi2012} applied to the precise
representative.
\end{proof}

\begin{lemma}[Differentiation of the distribution function]\label{lem:dist-fn}
Let $\Om\subset\R^2$ be bounded and let $u\in C^1(\Om)\cap C(\overline\Om)$
with $u\ge 0$. Set $m(t):=|\{u>t\}|$. Then $m$ is non-increasing, and for
Lebesgue-almost every $t$ with $\HH^1(\{u=t\}\cap\{\nabla u\ne 0\})=\HH^1(\{u=t\})$
one has
\[
  -m'(t)=\int_{\{u=t\}}\frac{1}{|\nabla u|}\,d\HH^1 .
\]
\end{lemma}

\begin{proof}
For $t<s$ we have, by the coarea formula (Theorem~\ref{thm:coarea}) applied
with $g=\mathbf 1_{\{t<u\le s\}}\,|\nabla u|^{-1}\mathbf 1_{\{\nabla u\ne0\}}$,
\[
  \int_{\{t<u\le s\}\cap\{\nabla u\ne0\}}1\,dx
  =\int_t^s\Big(\int_{\{u=\tau\}\cap\{\nabla u\ne0\}}\frac{1}{|\nabla u|}\,
  d\HH^1\Big)\,d\tau .
\]
By Sard's theorem (Theorem~\ref{thm:sard} below, valid since $u\in C^1$ and
$\R$ is $1$-dimensional, or directly because the critical values form a null
set) the set $\{\nabla u=0\}$ contributes a $\HH^1$-null portion of
$\{u=\tau\}$ for a.e.\ $\tau$, and $\{u=\tau\}\cap\{\nabla u=0\}$ has measure
zero in $\R^2$, hence
\[
  m(t)-m(s)=|\{t<u\le s\}|
  =\int_t^s\Big(\int_{\{u=\tau\}}\frac{1}{|\nabla u|}\,d\HH^1\Big)\,d\tau .
\]
Thus the non-increasing function $t\mapsto m(t)$ has, as its
(absolutely continuous part of the) distributional derivative, the function
$-\int_{\{u=t\}}|\nabla u|^{-1}\,d\HH^1$. Since $m$ is monotone it is
differentiable a.e., and at a.e.\ Lebesgue point $t$ of the right-hand
integrand the displayed integral identity yields
$-m'(t)=\int_{\{u=t\}}|\nabla u|^{-1}\,d\HH^1$.
\end{proof}

\subsection{Isoperimetric and relative isoperimetric inequalities}
\label{ss:isoperimetric}

\begin{theorem}[Planar isoperimetric inequality]\label{thm:isoperimetric}
For every set of finite perimeter $E\subset\R^2$ with $|E|<\infty$,
\[
  P(E)^2\ge 4\pi\,|E|,
\]
and equality holds if and only if $E$ is equivalent (up to a Lebesgue-null
set) to a disk $B_r(x)$.
\end{theorem}

\begin{proof}
See \cite[Theorem~14.1]{Maggi2012} for the inequality in $\R^n$ in the form
$P(E)\ge n\,\omega_n^{1/n}|E|^{(n-1)/n}$, which for $n=2$ reads
$P(E)\ge 2\sqrt{\pi}\,|E|^{1/2}$, i.e.\ $P(E)^2\ge 4\pi|E|$; the rigidity
statement (equality only for balls) is \cite[Theorem~14.1]{Maggi2012} as
well. See also \cite[Section~5.6.2, Theorem~5.6]{EvansGariepy2015}.
\end{proof}

\begin{theorem}[Relative isoperimetric inequality in a disk]
\label{thm:relative-isoperimetric}
There is an absolute constant $C>0$ such that for every set of finite
perimeter $E\subset\R^2$, every $x\in\R^2$ and every $r>0$,
\[
  \min\big\{\,|E\cap B_r(x)|,\ |B_r(x)\setminus E|\,\big\}
  \le C\,P\big(E;B_r(x)\big)^2 .
\]
The sharp constant is $C=1/\pi$.
\end{theorem}

\begin{proof}
This is the relative isoperimetric inequality on a convex domain; see
\cite[Theorem~2, Section~5.6.2]{EvansGariepy2015} for the existence of a
dimensional constant, and \cite[Section~12.5 and Theorem~14.4]{Maggi2012} for
the convex/relative formulation. The sharp constant $C=1/\pi$ in the plane
(attained as $E\cap B_r(x)$ degenerates to a half-disk through the centre,
where $|E\cap B_r|\to\frac12\pi r^2$ and $P(E;B_r)\to 2r$, giving the ratio
$\frac{\pi r^2/2}{(2r)^2}=\frac{\pi}{8}\le\frac1\pi$, the extremal half-disk
configuration $\min=\frac12\pi r^2$, $P=2r$ saturating $\min=\frac1\pi P^2$ at
$P=\sqrt{\pi/2}\,2r$) is the planar case of
\cite[Theorem~14.4]{Maggi2012}. For the purposes of this paper only the
existence of an absolute $C$ is used.
\end{proof}

\subsection{The sharp quantitative isoperimetric inequality}
\label{ss:fmp}

For a set of finite perimeter $E\subset\R^n$ with $0<|E|<\infty$, let $B(x)$
denote the ball centred at $x$ with $|B(x)|=|E|$, and define the
\emph{Fraenkel asymmetry}
\[
  \mathcal A(E):=\min_{x\in\R^n}\frac{|E\,\triangle\,B(x)|}{|E|},
\]
where $\triangle$ is symmetric difference, and the \emph{isoperimetric
deficit}
\[
  \delta(E):=\frac{P(E)}{P(B)}-1,
\]
with $B$ any ball of volume $|E|$ (so $P(B)=n\,\omega_n^{1/n}|E|^{(n-1)/n}$).
Both $\mathcal A(E)$ and $\delta(E)$ are scale- and translation-invariant and
nonnegative, vanishing precisely when $E$ is a ball.

\begin{theorem}[Fusco--Maggi--Pratelli]\label{thm:fmp}
There is a constant $C(n)>0$, depending only on the dimension $n$, such that
for every set of finite perimeter $E\subset\R^n$ with $0<|E|<\infty$,
\[
  \mathcal A(E)\le C(n)\,\sqrt{\delta(E)} .
\]
In particular, for $n=2$ there is an absolute constant $C(2)$ with
$\mathcal A(E)\le C(2)\sqrt{\delta(E)}$ for every planar set of finite
perimeter of positive finite area.
\end{theorem}

\begin{proof}
This is the main theorem of N.~Fusco, F.~Maggi and A.~Pratelli,
\emph{The sharp quantitative isoperimetric inequality}, Annals of Mathematics
\textbf{168} (2008), no.~3, 941--980, stated there as Theorem~1.1
\cite{FMP2008}. The power $\tfrac12$ on $\delta(E)$ is optimal (sharp): it
cannot be replaced by any larger exponent, as shown by ellipsoidal
perturbations of the ball.
\end{proof}

It is convenient to record the additive form of the planar inequality in
terms of the \emph{isoperimetric deficit by difference}
\[
  D_I(E):=P(E)^2-4\pi\,|E|\ \ge\ 0 ,
\]
which is nonnegative by Theorem~\ref{thm:isoperimetric}. The following lemma
is pure algebra given the FMP statement.

\begin{lemma}[Equivalence of additive and multiplicative deficits, $n=2$]
\label{lem:deficit-equiv}
Let $E\subset\R^2$ be a set of finite perimeter with $0<|E|<\infty$. Then
\[
  \delta(E)=\sqrt{\,1+\frac{D_I(E)}{4\pi|E|}\,}-1 .
\]
Consequently, there is an absolute constant $C>0$ such that whenever
$D_I(E)\le |E|$ one has
\[
  \mathcal A(E)^2\le C\,\frac{D_I(E)}{|E|} .
\]
\end{lemma}

\begin{proof}
In $\R^2$ a ball $B$ with $|B|=|E|$ satisfies $P(B)=2\sqrt{\pi|E|}$, hence
$P(B)^2=4\pi|E|$. Therefore
\[
  \big(1+\delta(E)\big)^2=\frac{P(E)^2}{P(B)^2}
  =\frac{P(E)^2}{4\pi|E|}
  =\frac{4\pi|E|+D_I(E)}{4\pi|E|}
  =1+\frac{D_I(E)}{4\pi|E|} .
\]
Since $1+\delta(E)>0$, taking the positive square root gives
$1+\delta(E)=\sqrt{1+D_I(E)/(4\pi|E|)}$, which is the asserted identity.

For the second statement, set $s:=D_I(E)/(4\pi|E|)\ge 0$. Then
$\delta(E)=\sqrt{1+s}-1=\dfrac{s}{\sqrt{1+s}+1}\le \dfrac{s}{2}$ for all
$s\ge0$, because $\sqrt{1+s}+1\ge 2$. Hence
\[
  \delta(E)\le\frac{s}{2}=\frac{D_I(E)}{8\pi|E|} .
\]
By Theorem~\ref{thm:fmp} with $n=2$,
$\mathcal A(E)^2\le C(2)^2\,\delta(E)\le \dfrac{C(2)^2}{8\pi}\,
\dfrac{D_I(E)}{|E|}$, which is the claim with
$C=C(2)^2/(8\pi)$. (The hypothesis $D_I(E)\le|E|$ is not needed for this
bound; it merely guarantees that $s$, and hence the regime, stays bounded so
that $\mathcal A(E)\le 2$ trivially and the estimate is non-vacuous.)
\end{proof}

\begin{remark}\label{rem:deficit-comparison}
The inequality $\delta(E)\le D_I(E)/(8\pi|E|)$ from
Lemma~\ref{lem:deficit-equiv} is one-sided in the sense used above; in the
reverse direction $\delta(E)\ge \dfrac{D_I(E)}{8\pi|E|(1+\delta(E))}$, so the
two deficits are comparable, $\delta(E)\asymp D_I(E)/|E|$, on any regime where
$\delta(E)$ is bounded. We will only use the upper bound.
\end{remark}

\subsection{Interior regularity and analyticity of the torsion function}
\label{ss:torsion}

Let $u\in H^1_0(\Om)$ be the weak solution of the torsion problem
\begin{equation}\label{eq:torsion}
  -\Delta u=1\ \text{ in }\Om,\qquad u=0\ \text{ on }\partial\Om,
\end{equation}
i.e.\ $\int_\Om\nabla u\cdot\nabla\varphi\,dx=\int_\Om\varphi\,dx$ for all
$\varphi\in H^1_0(\Om)$. Existence and uniqueness follow from the
Lax--Milgram theorem. We record its regularity.

\begin{theorem}[Smoothness and analyticity of $u$]\label{thm:torsion-regularity}
Let $u\in H^1_0(\Om)$ solve \eqref{eq:torsion} weakly. Then
$u\in C^\infty(\Om)$, and in fact $u$ is real-analytic in $\Om$.
\end{theorem}

\begin{proof}
The interior Schauder theory gives $u\in C^\infty(\Om)$: since the right-hand
side $f\equiv 1$ is smooth and $-\Delta$ has smooth (constant) coefficients,
interior Schauder estimates and a bootstrap argument yield
$u\in C^{k,\alpha}(\Om)$ for every $k$, hence $u\in C^\infty(\Om)$; see
D.~Gilbarg and N.~Trudinger, \emph{Elliptic Partial Differential Equations
of Second Order}, Springer, 2nd ed.\ 1983 (reprint of the 1998 edition),
Theorem~6.17 and Corollary~6.18 \cite{GT1983}. Real-analyticity in
the interior follows from the analyticity theory for solutions of elliptic
equations with analytic coefficients and analytic data: the Laplacian has
(constant, hence) real-analytic coefficients and the datum $1$ is analytic, so
$u$ is real-analytic in $\Om$ by C.~B.~Morrey, \emph{Multiple Integrals in
the Calculus of Variations}, Springer, 1966, Theorem~6.7.6 \cite{Morrey1966}.
(Equivalently, one may invoke Morrey--Nirenberg analytic regularity for linear
elliptic equations.)
\end{proof}

\begin{theorem}[Interior gradient and Hessian estimates]
\label{thm:interior-estimates}
There is an absolute constant $C>0$ with the following property. Let $u$ solve
\eqref{eq:torsion} in $\Om$, let $x\in\Om$, and set
$d:=\operatorname{dist}(x,\partial\Om)$. Then
\[
  |\nabla u(x)|\le C\Big(\frac{\|u\|_{L^\infty(B_d(x))}}{d}+d\Big),
  \qquad
  |D^2 u(x)|\le C\Big(\frac{\|u\|_{L^\infty(B_d(x))}}{d^2}+1\Big).
\]
\end{theorem}

\begin{proof}
Apply the interior Schauder estimate
\cite[Theorem~6.2, see also Corollary~6.3]{GT1983} (equivalently the
a priori interior estimate of \cite[Theorem~3.9]{GT1983} combined
with Schauder theory) to $u$ on the ball $B_d(x)$, where $-\Delta u=1$. The
scale-invariant form of these estimates gives, for the rescaled function,
$d\,\|\nabla u\|_{L^\infty(B_{d/2}(x))}
 +d^2\,\|D^2u\|_{L^\infty(B_{d/2}(x))}
 \le C\big(\|u\|_{L^\infty(B_d(x))}+d^2\|1\|_{L^\infty(B_d(x))}\big)$,
with $C$ absolute (depending only on dimension $n=2$ and the H\"older exponent
chosen for the datum, which here is irrelevant since $f\equiv1$). Dividing by
$d$ and by $d^2$ respectively and evaluating at $x$ yields the two stated
bounds.
\end{proof}

\begin{lemma}[A priori $L^\infty$ bound]\label{lem:linfty}
Let $u\in H^1_0(\Om)$ solve \eqref{eq:torsion} weakly, with $\Om\subset\R^2$
bounded. Then
\[
  0\le u(x)\le \tfrac14(\operatorname{diam}\Om)^2\qquad\text{for all }x\in\Om .
\]
\end{lemma}

\begin{proof}
By the weak maximum principle for $-\Delta u=1\ge0$ with $u=0$ on
$\partial\Om$ we have $u\ge0$ in $\Om$ (see
\cite[Theorem~8.1]{GT1983}); since $u\in C^\infty(\Om)\cap
C(\overline\Om)$ by Theorem~\ref{thm:torsion-regularity} and standard boundary
behaviour, this is a pointwise statement on $\Om$. For the upper bound, fix
$x_0$ such that $\Om\subset B_R(x_0)$ with $R=\operatorname{diam}\Om$ (any
$x_0\in\overline\Om$ works since every point of $\Om$ is within distance
$\operatorname{diam}\Om$ of $x_0$). Define the comparison function
\[
  w(x):=\frac{R^2-|x-x_0|^2}{4},
\]
which satisfies $-\Delta w=1$ in $\R^2$ and $w\ge0$ on $\overline\Om$ (indeed
$w\ge0$ on $B_R(x_0)\supseteq\Om$). Then $-\Delta(w-u)=0$ in $\Om$ and
$w-u=w\ge0$ on $\partial\Om$, so by the maximum principle
\cite[Corollary~3.2]{GT1983} $w-u\ge0$ in $\Om$, i.e.
$u\le w\le R^2/4=\tfrac14(\operatorname{diam}\Om)^2$.
\end{proof}

\subsection{Sard's theorem}\label{ss:sard}

\begin{theorem}[Sard]\label{thm:sard}
Let $U\subseteq\R^d$ be open and let $f\in C^k(U;\R^m)$ with
$k\ge\max\{1,d-m+1\}$. Then the set of \emph{critical values} of $f$,
\[
  f\big(\{x\in U:\operatorname{rank} Df(x)<m\}\big),
\]
has Lebesgue measure zero in $\R^m$. In particular, if $f\in C^\infty(U)$
with $m=1$ (so $k\ge d$ is automatic for $f\in C^\infty$), the set of
critical values $f(\{\nabla f=0\})\subset\R$ has Lebesgue measure zero.
\end{theorem}

\begin{proof}
This is Sard's theorem; see J.~Milnor, \emph{Topology from the Differentiable
Viewpoint}, Princeton University Press, revised ed.\ 1997, the chapter
``Sard's theorem'' (pp.~10--19), or the original A.~Sard, \emph{The measure of
the critical values of differentiable maps}, Bull.\ Amer.\ Math.\ Soc.\
\textbf{48} (1942), 883--890 \cite{Milnor1997,Sard1942}.
\end{proof}

\begin{lemma}[Almost every level of $u$ is regular]\label{lem:regular-values}
Let $u\in H^1_0(\Om)$ solve \eqref{eq:torsion}. Then for Lebesgue-almost every
$t\in\R$ the value $t$ is a regular value of $u$ in $\Om$, that is,
$\nabla u\ne0$ at every point of the level set $\{x\in\Om:u(x)=t\}$.
Consequently, for a.e.\ $t$ the level set $\{u=t\}\cap\Om$ is a (possibly
empty) smooth $1$-dimensional submanifold of $\Om$, and on it the coarea
quotient $1/|\nabla u|$ in Lemma~\ref{lem:dist-fn} is finite and smooth.
\end{lemma}

\begin{proof}
By Theorem~\ref{thm:torsion-regularity}, $u\in C^\infty(\Om)$ (indeed
real-analytic). Apply Sard's theorem (Theorem~\ref{thm:sard}) with $d=2$,
$m=1$: the set of critical values $u(\{x\in\Om:\nabla u(x)=0\})$ has Lebesgue
measure zero in $\R$. Hence for a.e.\ $t$, no point of $\{u=t\}\cap\Om$ is a
critical point, i.e.\ $\nabla u\ne0$ on the level set. For such regular $t$,
the implicit function theorem makes $\{u=t\}\cap\Om$ a smooth embedded
$1$-manifold, and $1/|\nabla u|$ is finite and smooth there.
\end{proof}


\bigskip
\section{Smoothness of the level sets}\label{sec:smooth}

Throughout this section $\Om\subset\R^2$ is a bounded open set that is
\emph{regular for the Dirichlet problem}, and $u$ denotes its torsion
function: the unique
\[
  u\in H^1_0(\Om)\cap C(\overline\Om),\qquad -\Delta u=1\ \text{in }\Om,
  \qquad u=0\ \text{on }\partial\Om .
\]
We collect from the Preliminaries (Section~\ref{sec:prelim}) the four facts
we are allowed to use without reproof:
\begin{enumerate}
\item[(P1)] $u$ is real-analytic in $\Om$ (interior analyticity of solutions
  of $-\Delta u=1$; see \cite[Ch.~6, \S6.1 and the analyticity
  theorem]{Morrey1966} or \cite[Theorem~6.6.1]{Morrey1966}).
\item[(P2)] Interior gradient and Hessian estimates: for every $K\Subset\Om$
  there is $C(K)$ with $\|u\|_{C^2(K)}\le C(K)$
  (\cite[Theorems~6.2, 6.17 and Corollary~6.3]{GT1983}).
\item[(P3)] Sard's theorem: the set of critical values of a $C^k$ map
  $\R^2\to\R$ with $k\ge2$ has Lebesgue measure zero
  (\cite[\S2--3]{Milnor1997}; \cite{Sard1942}).
\item[(P4)] The coarea formula and the structure of sets of finite perimeter
  (Section~\ref{sec:prelim}; \cite[\S3.4, \S5]{EvansGariepy2015}).
\end{enumerate}

\medskip
\noindent\textbf{Positivity of $u$.}
Since $-\Delta u=1>0$ in $\Om$, $u$ is superharmonic and not constant; as
$u=0$ on $\partial\Om$ and $u\not\equiv0$, the strong maximum principle
(equivalently Hopf's lemma) forbids an interior minimum equal to the boundary
value unless $u$ is constant, hence $u>0$ throughout the connected
components of $\Om$ meeting the interior; concretely, if $u(x_0)=0=\min u$ at
some $x_0\in\Om$ then the strong maximum principle for $-\Delta u\ge0$ would
force $u\equiv0$ near $x_0$, contradicting $-\Delta u=1$. Therefore $u>0$ in
$\Om$. Set
\[
  m:=\max_{\overline\Om}u=\max_{\Om}u>0 ,
\]
the maximum being attained by continuity on the compact set $\overline\Om$
and being positive since $u>0$ somewhere. For $v\in(0,m)$ write
\[
  E_v:=\{x\in\Om:u(x)>v\} .
\]

\subsection*{Compact containment of superlevel sets}

\begin{lemma}[Compact containment]\label{lem:compact-containment}
For every $v\in(0,m)$ the superlevel set
$\overline{E_v}=\{x\in\overline\Om:u(x)\ge v\}=\{u\ge v\}$ is a nonempty
compact subset of the open set $\Om$, and
\[
  \operatorname{dist}\bigl(\{u\ge v\},\partial\Om\bigr)\ \ge\ c(v)>0 .
\]
\end{lemma}

\begin{proof}
Since $u\in C(\overline\Om)$ and $\overline\Om$ is compact, the function
$u\colon\overline\Om\to\R$ is continuous on a compact set.

\emph{Nonempty.} By definition of $m$ there is $x_*\in\Om$ with $u(x_*)=m>v$,
so $x_*\in\{u\ge v\}$.

\emph{Closed and bounded.} The set $\{u\ge v\}=u^{-1}([v,\infty))$ is the
preimage of a closed set under the continuous map $u$ on $\overline\Om$, hence
relatively closed in $\overline\Om$; as $\overline\Om$ is itself closed in
$\R^2$, $\{u\ge v\}$ is closed in $\R^2$. It is contained in the bounded set
$\overline\Om$, hence bounded. A closed and bounded subset of $\R^2$ is
compact (Heine--Borel).

\emph{Disjoint from $\partial\Om$.} On $\partial\Om$ we have $u=0<v$, so no
boundary point lies in $\{u\ge v\}$. Thus $\{u\ge v\}\subset\Om$.

\emph{Positive distance.} The two sets $\{u\ge v\}$ and $\partial\Om$ are
disjoint and both compact (the first by the above, the second because
$\partial\Om$ is closed and bounded). The distance function
$x\mapsto\operatorname{dist}(x,\partial\Om)$ is continuous and strictly
positive on the compact set $\{u\ge v\}$ (strictly positive because that set
is disjoint from the closed set $\partial\Om$), so it attains a positive
minimum
\[
  c(v):=\min_{x\in\{u\ge v\}}\operatorname{dist}(x,\partial\Om)>0 .
\]
This is the asserted lower bound. Finally
$\overline{E_v}=\{u\ge v\}$: clearly $E_v\subset\{u\ge v\}$ and the latter is
closed, so $\overline{E_v}\subset\{u\ge v\}$; conversely if $u(x)\ge v$ then,
$u$ being continuous and nonconstant near $x$ (analytic, $-\Delta u=1$), every
neighbourhood of $x$ meets $\{u>v\}$, whence $x\in\overline{E_v}$. (We will
re-derive $\partial E_v=\{u=v\}$ more precisely in
Lemma~\ref{lem:smooth-level} for regular $v$.)
\end{proof}

\begin{remark}\label{rem:uniform-c2}
Combining Lemma~\ref{lem:compact-containment} with (P2): fixing
$v\in(0,m)$, the compact set $K_v:=\{u\ge v/2\}\Subset\Om$ contains an open
neighbourhood of $\{u\ge v\}$, and on $K_v$ all estimates (P1)--(P2) apply.
In particular $u\in C^\infty(K_v)$, indeed real-analytic, with $\nabla u$ and
$D^2u$ bounded there. This is what lets us apply the divergence theorem and
the implicit function theorem in a fixed interior region.
\end{remark}

\subsection*{Regular level sets are analytic curves}

Recall that $v\in(0,m)$ is a \emph{regular value} of $u$ if $\nabla u(x)\ne0$
for every $x\in\{u=v\}$; otherwise $v$ is a \emph{critical value}.

\begin{lemma}[Smooth regular level sets]\label{lem:smooth-level}
Let $v\in(0,m)$ be a regular value of $u$. Then:
\begin{enumerate}
\item[(a)] $\{u=v\}$ is a nonempty compact subset of $\Om$ on which
  $|\nabla u|>0$ everywhere;
\item[(b)] $\{u=v\}$ is a real-analytic embedded $1$-manifold without
  boundary, equal to a finite disjoint union of analytic Jordan curves
  $\Gamma_1,\dots,\Gamma_N$;
\item[(c)] $E_v=\{u>v\}$ is a bounded open subset of $\Om$ with topological
  boundary $\partial E_v=\{u=v\}$;
\item[(d)] $E_v$ has finite perimeter,
  $P(E_v)=\HH^1(\{u=v\})=\sum_{i=1}^N\HH^1(\Gamma_i)$, and its measure-theoretic
  outward unit normal is
  \[
    \nu(x)=-\frac{\nabla u(x)}{|\nabla u(x)|}\qquad\text{for every }x\in\{u=v\}.
  \]
\end{enumerate}
\end{lemma}

\begin{proof}
\emph{Step 0: $\{u=v\}$ is compact and nonempty.}
By Lemma~\ref{lem:compact-containment}, $\{u\ge v\}\Subset\Om$ is compact, and
$\{u=v\}=\{u\ge v\}\cap\{u\le v\}$ is closed (intersection of two closed sets),
hence compact as a closed subset of a compact set. It is nonempty: $u$ attains
both the value $m>v$ (at the interior maximum) and, on $\partial\Om$, the value
$0<v$; by the intermediate value theorem along any path in $\overline\Om$ from
the maximizer to a boundary point, $u$ takes the value $v$. (The path lies in
$\overline\Om$, $u$ is continuous on it, and $u$ runs from $m$ to $0$.)
Since $\{u=v\}\subset\{u\ge v\}\subset\Om$, it lies in the open set $\Om$, and
$|\nabla u|>0$ on it because $v$ is a regular value.

\emph{Step 1: local analytic graphs (implicit function theorem).}
Fix $p\in\{u=v\}$. By Remark~\ref{rem:uniform-c2}, $u$ is real-analytic in a
neighbourhood of $p$, and $\nabla u(p)\ne0$. After relabelling coordinates we
may assume $\partial_{x_2}u(p)\ne0$. The real-analytic implicit function
theorem (\cite[Theorem~6.1.2]{KrantzParks2002}; for the $C^1$ statement
\cite[\S0.3 and Appendix]{GT1983}) yields an open neighbourhood
$U_p\ni p$, an interval $I_p$, and a real-analytic function
$g_p\colon I_p\to\R$ such that
\[
  \{u=v\}\cap U_p=\{(x_1,g_p(x_1)):x_1\in I_p\}.
\]
Thus near $p$ the level set is the graph of an analytic function, hence an
embedded analytic $1$-dimensional submanifold; the map $x_1\mapsto(x_1,g_p(x_1))$
is an analytic chart. As $p$ was arbitrary, $\{u=v\}$ is a real-analytic
embedded $1$-manifold without boundary in $\Om$. (No boundary: every point has
a chart modelled on an open interval, with no half-line charts, because the
graph extends to both sides of $p$ inside $U_p$.)

\emph{Step 2: finitely many components.}
By Step~1 each $p\in\{u=v\}$ has an open neighbourhood $U_p$ meeting $\{u=v\}$
in a single connected arc. The family $\{U_p\}_{p\in\{u=v\}}$ is an open cover
of the compact set $\{u=v\}$, so admits a finite subcover
$U_{p_1},\dots,U_{p_M}$. A connected component $\Gamma$ of $\{u=v\}$ is open
in $\{u=v\}$ (locally it is an arc, hence locally connected, so components are
open) and closed (components are always closed), thus compact; being covered by
finitely many of the $U_{p_j}$, and each component being a union of arcs that
overlap, there can be at most $M$ components, i.e. finitely many. Call them
$\Gamma_1,\dots,\Gamma_N$, $N\le M$.

\emph{Step 3: each component is an analytic Jordan curve.}
Each $\Gamma_i$ is a connected, compact, $1$-dimensional analytic manifold
without boundary. By the classification of connected $1$-manifolds
(\cite[Appendix, Classification of $1$-manifolds]{Milnor1997}), a connected
compact $1$-manifold without boundary is diffeomorphic to the circle $S^1$;
moreover the diffeomorphism can be taken analytic by reparametrising by analytic
arc length (the manifold being analytic). Hence each $\Gamma_i$ is the image of
an analytic embedding $\gamma_i\colon S^1\hookrightarrow\Om$, i.e. an analytic
Jordan curve. (It cannot be a line: a noncompact model is excluded by the
compactness of $\Gamma_i$ from Step~2.)

\emph{Step 4: $E_v$ open, bounded, with $\partial E_v=\{u=v\}$.}
$E_v=u^{-1}((v,\infty))\cap\Om$ is the preimage of an open set under the
continuous map $u\colon\Om\to\R$, hence open; it is bounded, being contained in
$\overline\Om$. For the boundary, write $A:=\{u>v\}$, $B:=\{u<v\}\cup(\R^2\setminus\overline\Om)$.
Both $A$ and the interior of $B$ are open and disjoint from $\{u=v\}$, so
$\partial E_v\subset\{u=v\}$: indeed any $x\notin\{u=v\}$ either has $u(x)>v$
(then $x\in A$ open, interior point of $E_v$) or has $u(x)<v$ or $x\notin\overline\Om$
(then a neighbourhood avoids $E_v$, so $x$ is exterior). Conversely let
$x\in\{u=v\}$. By Step~1, in $U_x$ the set $\{u=v\}$ is the graph of $g_x$, and
since $\partial_{x_2}u(x)\ne0$, the function $x_2\mapsto u(x_1,x_2)$ is strictly
monotone across the graph; hence both $\{u>v\}$ and $\{u<v\}$ meet every
neighbourhood of $x$. Therefore $x\in\partial E_v$. Combining,
$\partial E_v=\{u=v\}$.

\emph{Step 5: finite perimeter, perimeter value, and normal.}
By Step~4, $\partial E_v=\{u=v\}$ is a compact $C^1$ (indeed analytic)
hypersurface in $\R^2$ contained in $\Om$, with $E_v$ lying locally on one side
(the side where $u>v$). A bounded open set with $C^1$ boundary is a set of
finite perimeter, and its reduced boundary coincides with its topological
boundary, with measure-theoretic outer normal equal to the classical outer
normal and
\[
  P(E_v)=\HH^1(\partial E_v)=\HH^1(\{u=v\})=\sum_{i=1}^N\HH^1(\Gamma_i)
\]
(\cite[\S5.7.2, Theorem~1 and \S5.1]{EvansGariepy2015}; this is also the
content of (P4)). It remains to identify the outward normal. The classical
outward unit normal of $E_v=\{u>v\}$ at $x\in\partial E_v$ points in the
direction of \emph{decreasing} $u$, i.e. opposite to $\nabla u$, since $\nabla u$
points in the direction of steepest increase of $u$. Normalising,
\[
  \nu(x)=-\frac{\nabla u(x)}{|\nabla u(x)|},
\]
which is well-defined because $|\nabla u(x)|>0$ on the regular level set. This
proves (a)--(d).
\end{proof}

\subsection*{The flux identity}

\begin{lemma}[Flux identity]\label{lem:flux}
For every regular value $v\in(0,m)$, with $\mu(v):=|E_v|$,
\[
  \int_{\{u=v\}}|\nabla u|\,d\HH^1=|E_v|=\mu(v),
  \qquad
  \int_{\{u=v\}}\frac{1}{|\nabla u|}\,d\HH^1=-\mu'(v),
\]
the second identity holding for a.e.\ regular $v$ (at which $\mu$ is
differentiable).
\end{lemma}

\begin{proof}
\emph{First identity.}
Fix a regular value $v$. By Lemma~\ref{lem:smooth-level}, $\overline{E_v}\Subset\Om$
is compact with smooth (analytic) boundary $\partial E_v=\{u=v\}$, and by
Remark~\ref{rem:uniform-c2} $u\in C^\infty$ in an open neighbourhood of
$\overline{E_v}$. The vector field $X:=\nabla u$ is therefore $C^1$ up to
$\overline{E_v}$, and the divergence theorem
(\cite[\S4 and the Gauss--Green theorem]{EvansGariepy2015};
\cite[(2.10)--(2.11)]{GT1983}) applies on the bounded $C^1$
domain $E_v$:
\[
  \int_{E_v}\operatorname{div}(\nabla u)\,dx
  =\int_{\partial E_v}\nabla u\cdot\nu\,d\HH^1 .
\]
On the left, $\operatorname{div}(\nabla u)=\Delta u=-1$, so the left side equals
$-|E_v|$. On the right, $\nu=-\nabla u/|\nabla u|$ by
Lemma~\ref{lem:smooth-level}(d), hence
\[
  \nabla u\cdot\nu=\nabla u\cdot\Bigl(-\frac{\nabla u}{|\nabla u|}\Bigr)
  =-\frac{|\nabla u|^2}{|\nabla u|}=-|\nabla u|,
\]
that is $\partial_\nu u=-|\nabla u|$ on $\{u=v\}$. Therefore
\[
  -|E_v|=\int_{\{u=v\}}\bigl(-|\nabla u|\bigr)\,d\HH^1,
\]
and multiplying by $-1$ gives
$\int_{\{u=v\}}|\nabla u|\,d\HH^1=|E_v|=\mu(v)$, the first identity.

\emph{Second identity.}
We use the coarea formula (P4): for the locally Lipschitz function $u$ on the
open set $\{u>0\}$ and any Borel $g\ge0$,
\[
  \int_{\{u>0\}}g(x)\,|\nabla u(x)|\,dx
  =\int_0^{\infty}\Bigl(\int_{\{u=t\}}g\,d\HH^1\Bigr)dt
  \qquad(\cite[\S3.4.2, Theorem~5]{EvansGariepy2015}).
\]
Apply this with $g=\mathbf 1_{\{v<u\le m\}}\,|\nabla u|^{-1}$ on the set
$\{0<u\}$, noting $|\nabla u|^{-1}$ is finite $\HH^1$-a.e.\ on $\{u=t\}$ for
a.e.\ regular $t$ and that $g|\nabla u|=\mathbf 1_{\{v<u\le m\}}$. Equivalently,
choosing $g=\mathbf 1_{\{u>v\}}/|\nabla u|$ where defined, the left side is
$|\{u>v\}|=\mu(v)$ minus boundary contributions of measure zero; the cleanest
route is to differentiate the coarea representation of $\mu$. Indeed, taking
$g=\mathbf 1_{\{u>v\}}$ in the coarea formula written for the measure of
superlevel sets gives, for $0<v<m$,
\[
  \mu(v)=|\{u>v\}|=\int_{\{u>v\}}\!1\,dx
  =\int_{\{u>v\}}\frac{|\nabla u|}{|\nabla u|}\,dx
  =\int_v^{m}\Bigl(\int_{\{u=t\}}\frac{1}{|\nabla u|}\,d\HH^1\Bigr)dt,
\]
where in the last step we applied coarea with $g=\mathbf 1_{\{u>v\}}/|\nabla u|$
and used that $|\{\nabla u=0\}|=0$ (P3/Sard implies the critical \emph{values}
have measure zero; the critical \emph{set} $\{\nabla u=0\}$ has Lebesgue measure
zero because $u$ is real-analytic and nonconstant, so $\{\nabla u=0\}$ is a
proper analytic subvariety, hence $\HH^2$-null). Thus $\mu$ is, on $(0,m)$, the
integral from $v$ to $m$ of the function
\[
  v\mapsto \Phi(v):=\int_{\{u=v\}}\frac{1}{|\nabla u|}\,d\HH^1 ,
\]
which is finite for a.e.\ regular $v$. By the Lebesgue differentiation theorem
(fundamental theorem of calculus for the absolutely continuous function
$v\mapsto\int_v^m\Phi$), for a.e.\ $v$,
\[
  \mu'(v)=-\Phi(v)=-\int_{\{u=v\}}\frac{1}{|\nabla u|}\,d\HH^1 ,
\]
i.e. $\int_{\{u=v\}}|\nabla u|^{-1}\,d\HH^1=-\mu'(v)$, the second identity. The
minus sign reflects that $\mu$ is non-increasing (Lemma~\ref{lem:critvals}).
\end{proof}

\subsection*{Critical values and regularity of $\mu$}

\begin{lemma}[Finitely-thin critical set; regularity of $\mu$]\label{lem:critvals}
Let $\mathrm{Crit}:=\{v\in(0,m):v\text{ is a critical value of }u\}$. Then:
\begin{enumerate}
\item[(a)] $\mathrm{Crit}$ is a closed subset of $(0,m)$ of Lebesgue measure
  zero; consequently a.e.\ $v\in(0,m)$ is a regular value, and the conclusions of
  Lemmas~\ref{lem:smooth-level}--\ref{lem:flux} hold for a.e.\ $v$;
\item[(b)] for \emph{every} $v\in(0,m)$ the level set $\{u=v\}$ is Lebesgue-null:
  $|\{u=v\}|=\HH^2(\{u=v\})=0$;
\item[(c)] $\mu(v)=|E_v|$ is non-increasing on $(0,m)$ and continuous, with no
  jumps (no atoms).
\end{enumerate}
\end{lemma}

\begin{proof}
\emph{(a) Measure zero and closedness.}
The critical \emph{values} of $u$ have measure zero by Sard's theorem (P3),
applicable since $u\in C^2(\Om)$ (indeed $C^\infty$, by (P1)--(P2)). For
closedness: let $v_k\in\mathrm{Crit}$, $v_k\to v\in(0,m)$. Choose
$x_k\in\{u=v_k\}$ with $\nabla u(x_k)=0$. By
Lemma~\ref{lem:compact-containment}, all $x_k$ lie in the fixed compact set
$\{u\ge \tfrac{1}{2}\min_k v_k\}\Subset\Om$ (and $\min_k v_k>0$ since
$v_k\to v>0$); passing to a subsequence, $x_k\to x_*\in\Om$. By continuity of
$u$ and of $\nabla u$ (P2), $u(x_*)=v$ and $\nabla u(x_*)=0$, so $v$ is a
critical value, i.e. $v\in\mathrm{Crit}$. Hence $\mathrm{Crit}$ is closed in
$(0,m)$.

\emph{(b) Level sets are area-null.}
Two cases. If $v$ is a regular value, then by Lemma~\ref{lem:smooth-level}
$\{u=v\}$ is a finite union of analytic curves, each of which has zero
$2$-dimensional Lebesgue measure (a $1$-manifold in $\R^2$ is locally a graph,
hence $\HH^2$-null), so $|\{u=v\}|=0$. If $v$ is a critical value, write
\[
  \{u=v\}=\bigl(\{u=v\}\cap\{\nabla u\ne0\}\bigr)\cup\bigl(\{u=v\}\cap\{\nabla u=0\}\bigr).
\]
At points of the first set the implicit function theorem (as in
Lemma~\ref{lem:smooth-level}, Step~1) again represents $\{u=v\}$ locally as an
analytic graph, hence that part is $\HH^2$-null. The second set is contained in
$\{\nabla u=0\}$, which is Lebesgue-null because $u$ is real-analytic and
nonconstant on the connected components of $\Om$ (so $\nabla u$ does not vanish
identically): the zero set of a nonzero real-analytic function has Lebesgue
measure zero (\cite[Lemma in \S0; or the analytic zero-set theorem]{KrantzParks2002}).
Covering $\{u=v\}$ by countably many balls on which $u$ is analytic and applying
this in each, we get $|\{u=v\}|=0$ for every $v$.

\emph{(c) Monotonicity and continuity of $\mu$.}
Monotonicity: if $v_1<v_2$ then $E_{v_2}=\{u>v_2\}\subset\{u>v_1\}=E_{v_1}$, so
$\mu(v_2)=|E_{v_2}|\le|E_{v_1}|=\mu(v_1)$; thus $\mu$ is non-increasing.
Continuity: for $v_1<v_2$ in $(0,m)$,
\[
  \mu(v_1)-\mu(v_2)=|E_{v_1}\setminus E_{v_2}|=|\{v_1<u\le v_2\}|.
\]
As $v_2\downarrow v_1$, the sets $\{v_1<u\le v_2\}$ decrease to
$\{u=v_1\}\cup\varnothing=\{u=v_1\}$ (more precisely
$\bigcap_{v_2>v_1}\{v_1<u\le v_2\}=\{u=v_1\}$), which has measure $0$ by part
(b); by dominated convergence (all sets contained in the finite-measure set
$\Om$) the difference tends to $0$, giving right-continuity. Similarly, as
$v_1\uparrow v_2$ the sets $\{v_1<u<v_2\}$ decrease to $\varnothing$, while
$\{v_1<u\le v_2\}\downarrow\{u=v_2\}$, again of measure $0$, giving
left-continuity. Hence $\mu$ is continuous on $(0,m)$ with no atoms.
\end{proof}

\begin{remark}\label{rem:ae-suffices}
The combination of Lemmas~\ref{lem:smooth-level}--\ref{lem:critvals} is exactly
what later sections require: away from the measure-zero set $\mathrm{Crit}$,
the superlevel sets $E_v$ are smooth finite-perimeter domains for which the
flux identity and the coarea representation of $\mu$ hold pointwise, while the
exceptional set $\mathrm{Crit}$ may be discarded in any integration in $v$
(e.g. when integrating the flux identity against $dv$). When $u$ is moreover
real-analytic, $\mathrm{Crit}$ is in fact finite, but we shall only use that it
is closed and null.
\end{remark}


\bigskip
\section{The level-set deficit identity}\label{sec:identity}

We now exploit the torsion equation quantitatively. Throughout,
$\Om\subset\R^2$ is bounded open and regular for the Dirichlet problem,
$|\Om|=\pi$, $u$ is its torsion function (Section~\ref{sec:smooth}), and
\[
   T(\Om):=\int_\Om u\,dx=\int_\Om|\nabla u|^2\,dx,
   \qquad
   \dT:=\dT(\Om):=T(B_1)-T(\Om)=\frac{\pi}{8}-T(\Om)\ \ge\ 0 ,
\]
where $B_1$ is the unit disk, $T(B_1)=\pi/8$, and the inequality $\dT\ge0$
is the Saint--Venant inequality (a special case of
Theorem~\ref{thm:fmp}-type rigidity, or of Talenti's theorem below). The
two equalities defining $T(\Om)$ hold because, testing
$-\Delta u=1$ against $u$, $\int_\Om|\nabla u|^2=\int_\Om u$.

For a regular value $v\in(0,m)$ ($m=\max_\Om u$) we abbreviate, using the
notation of Section~\ref{sec:smooth},
\[
   \mu(v):=|E_v|,\qquad P(v):=P(E_v)=\HH^1(\{u=v\}).
\]
Lemma~\ref{lem:flux} gives the two structural identities
\begin{equation}\label{eq:flux-coarea}
   \int_{\{u=v\}}|\nabla u|\,d\HH^1=\mu(v),
   \qquad
   \int_{\{u=v\}}\frac{1}{|\nabla u|}\,d\HH^1=-\mu'(v)\quad(\text{a.e. }v).
\end{equation}

\begin{definition}[The two level deficits]\label{def:deficits}
For a regular value $v\in(0,m)$ set
\[
   D_H(v):=\mu(v)\,\bigl(-\mu'(v)\bigr)-P(v)^2,
   \qquad
   D_I(v):=P(v)^2-4\pi\,\mu(v).
\]
\end{definition}

\begin{lemma}[Non-negativity]\label{lem:nonneg}
For a.e.\ $v\in(0,m)$ one has $D_H(v)\ge0$ and $D_I(v)\ge0$. Moreover
$D_H(v)=0$ if and only if $|\nabla u|$ is constant $\HH^1$-a.e.\ on
$\{u=v\}$, and $D_I(v)=0$ if and only if $E_v$ is a disk.
\end{lemma}

\begin{proof}
$D_I(v)\ge0$ is the planar isoperimetric inequality
(Theorem~\ref{thm:isoperimetric}) applied to $E_v$, with equality iff $E_v$
is a disk. For $D_H$, the Cauchy--Schwarz inequality on the measure
$\HH^1\llcorner\{u=v\}$ gives
\[
   P(v)^2=\Bigl(\int_{\{u=v\}}1\,d\HH^1\Bigr)^2
   \le\Bigl(\int_{\{u=v\}}|\nabla u|\,d\HH^1\Bigr)
        \Bigl(\int_{\{u=v\}}\tfrac1{|\nabla u|}\,d\HH^1\Bigr)
   =\mu(v)\bigl(-\mu'(v)\bigr),
\]
by \eqref{eq:flux-coarea}; this is $D_H(v)\ge0$, with equality iff
$|\nabla u|$ and $1/|\nabla u|$ are proportional $\HH^1$-a.e.\ on
$\{u=v\}$, i.e.\ iff $|\nabla u|$ is $\HH^1$-a.e.\ constant there.
\end{proof}

\begin{theorem}[Level-set deficit identity]\label{thm:identity}
With the notation above,
\begin{equation}\label{eq:identity}
   \int_0^m\bigl(D_H(v)+D_I(v)\bigr)\,dv\ =\ 4\pi\,\dT(\Om).
\end{equation}
\end{theorem}

\begin{proof}
By Definition~\ref{def:deficits} the integrand telescopes: the $P(v)^2$
terms cancel and
\begin{equation}\label{eq:DHDI-sum}
   D_H(v)+D_I(v)=\mu(v)\bigl(-\mu'(v)\bigr)-4\pi\,\mu(v)
   =-\mu(v)\bigl(\mu'(v)+4\pi\bigr).
\end{equation}
We integrate the two terms separately over $(0,m)$. First, $\mu$ is
absolutely continuous on $(0,m)$: by the second identity in
\eqref{eq:flux-coarea} and Fubini,
$\mu(v)=\int_v^m\Phi(\tau)\,d\tau$ with
$\Phi(\tau)=\int_{\{u=\tau\}}|\nabla u|^{-1}d\HH^1\ge0$ and
$\int_0^m\Phi\,d\tau=\int_{\{0<u<m\}}1\,dx=|\Om|<\infty$ (coarea,
Theorem~\ref{thm:coarea}); hence $\Phi\in L^1(0,m)$ and $\mu$ is absolutely
continuous with $\mu'=-\Phi$ a.e. Therefore $v\mapsto\tfrac12\mu(v)^2$ is
absolutely continuous with derivative $\mu\mu'$, and
\[
   \int_0^m\bigl(-\mu(v)\mu'(v)\bigr)\,dv
   =-\Bigl[\tfrac12\mu(v)^2\Bigr]_0^m
   =\tfrac12\mu(0)^2-\tfrac12\mu(m)^2
   =\tfrac12|\Om|^2,
\]
since $\mu(0)=|\{u>0\}|=|\Om|$ (as $u>0$ in $\Om$) and
$\mu(m)=|\{u>m\}|=0$. Second, by the layer-cake formula,
\[
   \int_0^m 4\pi\,\mu(v)\,dv
   =4\pi\int_0^m|\{u>v\}|\,dv
   =4\pi\int_\Om u\,dx=4\pi\,T(\Om).
\]
Combining with \eqref{eq:DHDI-sum},
\[
   \int_0^m\bigl(D_H+D_I\bigr)\,dv
   =\tfrac12|\Om|^2-4\pi T(\Om)
   =\tfrac{\pi^2}{2}-4\pi T(\Om)
   =4\pi\Bigl(\tfrac{\pi}{8}-T(\Om)\Bigr)=4\pi\,\dT,
\]
using $|\Om|=\pi$ and $T(B_1)=\pi/8$.
\end{proof}

\begin{remark}[Interpretation]\label{rem:identity}
Identity \eqref{eq:identity} writes the Saint--Venant deficit as the
total, over all levels, of two non-negative geometric defects of the level
sets: the \emph{Serrin/H\"older defect} $D_H$, which measures the failure
of $|\nabla u|$ to be constant on $\{u=v\}$, and the \emph{isoperimetric
defect} $D_I$, which measures the failure of $E_v$ to be a disk. It is the
engine of everything that follows: if $\dT$ is small, then for most levels
both defects are small.
\end{remark}

\section{The volume profile: Talenti's comparison}\label{sec:talenti}

To turn ``most levels'' into a usable statement we need to know the
\emph{volume} of the level sets, i.e.\ the profile $\mu(v)$. This is
supplied by Talenti's pointwise comparison theorem.

\begin{theorem}[Talenti]\label{thm:talenti}
Let $u$ be the torsion function of $\Om$ with $|\Om|=\pi$, and let $u_B$ be
the torsion function of the disk $B_1$, namely
$u_B(x)=\tfrac14(1-|x|^2)$. Let $u^\ast$ denote the Schwarz (decreasing
radial) rearrangement of $u$. Then $u^\ast(x)\le u_B(x)$ for all $x\in
B_1$; equivalently, in terms of distribution functions,
\begin{equation}\label{eq:talenti}
   \mu(v)=|\{u>v\}|\ \le\ |\{u_B>v\}|=\pi\,(1-4v)_+\qquad(v\ge0).
\end{equation}
\end{theorem}

\begin{proof}
This is the model case of G.~Talenti, \emph{Elliptic equations and
rearrangements}, Ann.\ Scuola Norm.\ Sup.\ Pisa Cl.\ Sci.\ (4)
\textbf{3} (1976), 697--718, Theorem~1: for $-\Delta u=f$ in $\Om$ with
$u\in H^1_0(\Om)$, one has $u^\ast\le \tilde u$ pointwise, where $\tilde u$
solves $-\Delta\tilde u=f^\ast$ in the ball $\Om^\ast$ of the same volume
with zero boundary data. Here $f\equiv1$, $f^\ast\equiv1$, $\Om^\ast=B_1$,
and $\tilde u=u_B$. The distributional form \eqref{eq:talenti} follows
because Schwarz rearrangement preserves the distribution function,
$|\{u>v\}|=|\{u^\ast>v\}|\le|\{u_B>v\}|$, and
$\{u_B>v\}=\{|x|^2<1-4v\}$ has area $\pi(1-4v)_+$.
\end{proof}

\begin{corollary}[The maximum and the profile are pinned]\label{cor:profile}
There is an absolute constant $\delta_\star\in(0,\pi/16]$ such that if
$\dT(\Om)\le\delta_\star$ then
\begin{equation}\label{eq:m-pin}
   \frac{1}{16}\ \le\ m=\max_\Om u\ \le\ \frac14,
   \qquad
   \Bigl|m-\tfrac14\Bigr|\le\Bigl(\tfrac{\dT}{2\pi}\Bigr)^{1/2},
\end{equation}
and the profile is close to the ball profile in $L^1$:
\begin{equation}\label{eq:profile-L1}
   \int_0^m\Bigl(\pi(1-4v)_+-\mu(v)\Bigr)\,dv\ \le\ \dT .
\end{equation}
\end{corollary}

\begin{proof}
The upper bound $m\le\frac14$ is \eqref{eq:talenti} at $v=m$: if $m>\frac14$
then $\mu(m)\le\pi(1-4m)_+=0$ forces $\mu(m)=0$, consistent, but for
$v$ slightly below $m$, $\mu(v)>0$ while $\pi(1-4v)_+=0$, contradicting
\eqref{eq:talenti}; hence $m\le\frac14$. For the lower bound on $m$ and the
profile estimate, integrate \eqref{eq:talenti}:
\[
   T(\Om)=\int_0^m\mu(v)\,dv\ \le\ \int_0^{1/4}\pi(1-4v)\,dv=\frac{\pi}{8},
\]
recovering $\dT\ge0$, and more precisely, since
$\int_0^m\pi(1-4v)\,dv=\pi(m-2m^2)=\tfrac{\pi}{8}-2\pi\bigl(m-\tfrac14\bigr)^2$,
\[
   \dT=\frac{\pi}{8}-T(\Om)
   =\Bigl(\frac\pi8-\!\int_0^m\!\pi(1-4v)\,dv\Bigr)
     +\!\int_0^m\!\bigl(\pi(1-4v)-\mu(v)\bigr)\,dv
   =2\pi\bigl(m-\tfrac14\bigr)^2+R,
\]
where $R:=\int_0^m(\pi(1-4v)_+-\mu)\,dv\ge0$ by \eqref{eq:talenti}
(note $\pi(1-4v)_+=\pi(1-4v)$ on $(0,m)$ since $m\le\frac14$). As both
terms are non-negative and sum to $\dT$, each is $\le\dT$; this gives
\eqref{eq:profile-L1} and $|m-\frac14|\le(\dT/2\pi)^{1/2}$. Finally
$|m-\frac14|\le(\delta_\star/2\pi)^{1/2}\le\frac3{16}$ once
$\delta_\star\le 2\pi(3/16)^2$, so $m\ge\frac14-\frac3{16}=\frac1{16}$.
\end{proof}

\section{The good inset}\label{sec:goodlevel}

We now select one level set --- the \emph{inset} --- that is simultaneously
nearly isoperimetric, has nearly constant boundary gradient, and encloses a
definite area.

\begin{proposition}[Existence of a good inset]\label{prop:goodlevel}
There are absolute constants $\delta_\star>0$ and $C_0$ such that if
$\dT(\Om)\le\delta_\star$ there exists a regular value
$\hat v\in\bigl(\tfrac{m}{4},\tfrac{m}{2}\bigr)$ of $u$ with
\begin{equation}\label{eq:good}
   D_H(\hat v)+D_I(\hat v)\ \le\ C_0\,\dT,
   \qquad
   \mu(\hat v)\ \ge\ \frac\pi4 .
\end{equation}
We call $E_{\hat v}=\{u>\hat v\}$ the \emph{inset} and write
$\Etil:=E_{\hat v}$, $\hat\mu:=\mu(\hat v)\in[\tfrac\pi4,\pi)$,
$\hat P:=P(\hat v)$.
\end{proposition}

\begin{proof}
Work on the interval $J:=(\tfrac m4,\tfrac m2)$, of length $m/4\ge1/64$ by
Corollary~\ref{cor:profile}. By Theorem~\ref{thm:identity},
$\int_J(D_H+D_I)\,dv\le4\pi\dT$, so the set
\[
   J_{\rm bad}:=\Bigl\{v\in J:\ D_H(v)+D_I(v)>\tfrac{16\pi\dT}{m}\Bigr\}
\]
has measure $|J_{\rm bad}|<\dfrac{4\pi\dT}{16\pi\dT/m}=\dfrac m4=|J|$ by
Chebyshev's inequality; thus $|J\setminus J_{\rm bad}|>0$. On
$J\setminus J_{\rm bad}$, $D_H+D_I\le 16\pi\dT/m\le 256\pi\,\dT$ (using
$m\ge\tfrac1{16}$).

Next we control the volume. By \eqref{eq:talenti}, for $v\in J$ we have
$\pi(1-4v)\ge\pi(1-4\cdot\tfrac m2)=\pi(1-2m)\ge\pi(1-2\cdot\tfrac14)=\tfrac\pi2$.
Define
\[
   J_{\rm small}:=\Bigl\{v\in J:\ \mu(v)<\tfrac\pi4\Bigr\}
   \subseteq\Bigl\{v\in J:\ \pi(1-4v)-\mu(v)>\tfrac\pi4\Bigr\}.
\]
By Chebyshev applied to \eqref{eq:profile-L1},
$|J_{\rm small}|\le\dfrac{\int_0^m(\pi(1-4v)-\mu)\,dv}{\pi/4}\le\dfrac{4\dT}{\pi}$.
Hence
\[
   \bigl|J\setminus(J_{\rm bad}\cup J_{\rm small})\bigr|
   \ \ge\ |J|-|J_{\rm bad}|-|J_{\rm small}|
   \ >\ \frac m4-\frac m4 \ \cdots
\]
this last step needs $|J_{\rm bad}|+|J_{\rm small}|<|J|=m/4$. We have
$|J_{\rm bad}|<m/4$ but must leave room for $J_{\rm small}$; we instead use
a slightly smaller threshold. Replace $J_{\rm bad}$ by
$J_{\rm bad}':=\{v\in J:D_H+D_I>32\pi\dT/m\}$, so by Chebyshev
$|J_{\rm bad}'|<\tfrac{4\pi\dT}{32\pi\dT/m}=\tfrac m8$. Then
\[
   |J_{\rm bad}'|+|J_{\rm small}|<\frac m8+\frac{4\dT}\pi
   \ \le\ \frac m8+\frac{4\delta_\star}\pi\ <\ \frac m4=|J|
\]
provided $\delta_\star\le \pi m/32$, which holds for
$\delta_\star\le\pi/512$ (as $m\ge1/16$). Therefore
$J\setminus(J_{\rm bad}'\cup J_{\rm small})$ has positive measure;
since the critical values form a null set
(Lemma~\ref{lem:critvals}), we may pick a regular value $\hat v$ in it.
It satisfies $D_H(\hat v)+D_I(\hat v)\le 32\pi\dT/m\le512\pi\,\dT=:C_0\dT$
and $\mu(\hat v)\ge\pi/4$.
\end{proof}

\section{Controlled deficits of the inset}\label{sec:controlled}

We translate \eqref{eq:good} into the three geometric controls used in the
sequel: small isoperimetric deficit, small (weighted) Serrin deficit, and
$L^1$-closeness to a disk.

\begin{proposition}[Isoperimetric and Serrin control]\label{prop:controlled}
Let $\Etil=E_{\hat v}$ be the inset of Proposition~\ref{prop:goodlevel}.
Then, with $C_0$ as there and absolute constants,
\begin{align}
   \diso(\Etil):=\frac{\hat P}{2\sqrt{\pi\hat\mu}}-1
   &\ \le\ \frac{D_I(\hat v)}{8\pi\hat\mu}\ \le\ \frac{C_0}{2\pi^2}\,\dT,
   \label{eq:diso-inset}\\[1mm]
   W(\Etil):=\int_{\{u=\hat v\}}\frac{(|\nabla u|-c)^2}{|\nabla u|}\,d\HH^1
   &\ =\ \frac{\hat\mu}{\hat P^{2}}\,D_H(\hat v)\ \le\ \frac{C_0}{4\pi}\,\dT,
   \qquad c:=\frac{\hat\mu}{\hat P}=\frac{|\Etil|}{P(\Etil)}.
   \label{eq:serrin-inset}
\end{align}
\end{proposition}

\begin{proof}
For \eqref{eq:diso-inset}: $\hat P^2=4\pi\hat\mu+D_I(\hat v)$, so
$\bigl(1+\diso(\Etil)\bigr)^2=\hat P^2/(4\pi\hat\mu)=1+D_I(\hat
v)/(4\pi\hat\mu)$, whence (as in Lemma~\ref{lem:deficit-equiv})
$\diso(\Etil)\le D_I(\hat v)/(8\pi\hat\mu)$. Using
$\hat\mu\ge\pi/4$ and $D_I(\hat v)\le C_0\dT$ gives the bound.

For \eqref{eq:serrin-inset}: with $c=\hat\mu/\hat P$ and the identities
\eqref{eq:flux-coarea} (so $\int|\nabla u|=\hat\mu$,
$\int1/|\nabla u|=-\mu'(\hat v)$, $\int 1=\hat P$),
\[
   \int\frac{(|\nabla u|-c)^2}{|\nabla u|}
   =\int|\nabla u|-2c\hat P+c^2\!\int\frac1{|\nabla u|}
   =\hat\mu-2c\hat P+c^2(-\mu'(\hat v)).
\]
Substituting $c=\hat\mu/\hat P$ gives
$\hat\mu-2\hat\mu+\frac{\hat\mu^2}{\hat P^2}(-\mu'(\hat v))
=\frac{\hat\mu}{\hat P^2}\bigl(\hat\mu(-\mu'(\hat v))-\hat P^2\bigr)
=\frac{\hat\mu}{\hat P^2}D_H(\hat v)$.
Finally $\hat P^2\ge4\pi\hat\mu$ gives
$\hat\mu/\hat P^2\le1/(4\pi)$, and $D_H(\hat v)\le C_0\dT$.
\end{proof}

\begin{corollary}[$L^1$-closeness to a disk]\label{cor:L1}
There are a point $z\in\R^2$ and absolute constants $C,\delta_\star$ such
that if $\dT\le\delta_\star$, the inset satisfies
\begin{equation}\label{eq:L1}
   \bigl|\Etil\,\triangle\,B_{\hat r}(z)\bigr|\ \le\ C\,\hat\mu\,\sqrt{\dT},
   \qquad \hat r:=\sqrt{\hat\mu/\pi},
\end{equation}
where $B_{\hat r}(z)$ is the disk of area $\hat\mu=|\Etil|$ minimising the
symmetric difference.
\end{corollary}

\begin{proof}
By the sharp quantitative isoperimetric inequality
(Theorem~\ref{thm:fmp}, $n=2$) and Lemma~\ref{lem:deficit-equiv} applied to
$\Etil$,
\[
   \mathcal A(\Etil)=\min_z\frac{|\Etil\triangle B_{\hat r}(z)|}{\hat\mu}
   \le C\sqrt{\diso(\Etil)}\le C'\sqrt{\dT}
\]
by \eqref{eq:diso-inset}. Multiplying by $\hat\mu$ gives \eqref{eq:L1}.
\end{proof}

\begin{remark}[Normalisation]\label{rem:normalise}
Rescaling $\Etil$ by the factor $s:=\sqrt{\pi/\hat\mu}\in(1,2]$ (since
$\hat\mu\in[\pi/4,\pi)$) produces a set of area $\pi$ to which all the above
apply verbatim: $\diso$, $\mathcal A$ and the relative quantities below are
scale-invariant, while perimeters and lengths scale by $s$ and areas by
$s^2$. We henceforth state the geometric conclusions for the unrescaled
$\Etil$ and freely pass to the area-$\pi$ normalisation when convenient.
\end{remark}

\section{Topological regularity: no interior holes}
\label{sec:topology}

A consequence of the superharmonicity of $u$ ($-\Delta u=1>0$, so
$\Delta u<0$) and the minimum principle. The inner ball and the trapping
between two concentric circles are deferred to Section~\ref{sec:trapping},
where they follow from the smoothness of the inset.

\begin{proposition}[No interior holes]\label{prop:noholes}
For every $v\in(0,m)$, each connected component of the open set
$\{x\in\Om:u(x)<v\}$ has its closure meeting $\partial\Om$. Equivalently,
$\Om\setminus\overline{E_v}$ has no connected component compactly contained
in $\Om$: the inset $\Etil=E_{\hat v}$ has no interior holes.
\end{proposition}

\begin{proof}
Suppose some component $K$ of $\{u<v\}$ satisfied $\overline K\Subset\Om$.
Then $\overline K$ is compact in $\Om$, $u<v$ in $K$, and $u=v$ on
$\partial K$ (because $\partial K\subset\{u=v\}$: a boundary point of a
component of $\{u<v\}$ lying in $\Om$ cannot have $u<v$, as that is an open
condition placing it in the interior of $K$, nor $u>v$, as $K$ accumulates
there). Thus $u$ attains on the compact set $\overline K$ an interior
minimum value $<v=\min_{\partial K}u$. But $-\Delta u=1\ge0$ makes $u$
superharmonic, so it satisfies the minimum principle
\cite[Theorem~3.1 and the strong form, Theorem~3.5]{GT1983}:
$\min_{\overline K}u=\min_{\partial K}u$. This contradiction shows no such
$K$ exists.
\end{proof}


\bigskip
\section{The two-sided annulus via Magnanini--Poggesi}\label{sec:trapping}

We come to the ``squeezing between two radii.'' Here the smoothness of the
inset is decisive: the inset is a \emph{smooth} domain on which the boundary
gradient $|\nabla u|$ is nearly constant
(Proposition~\ref{prop:controlled}), i.e.\ it is a near-extremal domain for
\emph{Serrin's overdetermined problem}. The quantitative stability theory
for that problem applies to smooth domains and traps the inset between two
concentric circles. We emphasise at the outset (Remark~\ref{rem:nonuniform})
that the resulting width is controlled by a constant depending on the
inset's $C^2$ geometry, which is finite (the inset is analytic) but is not,
in this soft theory, uniform in $\dT$.

We first pass from the weighted Serrin deficit \eqref{eq:serrin-inset} to
the unweighted one used in the literature.

\begin{lemma}[Unweighted Serrin deficit]\label{lem:unweighted}
Let $\Etil=E_{\hat v}$ be the inset and
$G_1:=\max_{\{u=\hat v\}}|\nabla u|<\infty$ (finite, since $\{u=\hat v\}$ is
compact in $\Om$ and $u\in C^\infty$ near it). Then, with $c=\hat\mu/\hat P$,
\begin{equation}\label{eq:unweighted}
   \eps(\Etil)^2:=\int_{\{u=\hat v\}}\bigl(|\nabla u|-c\bigr)^2\,d\HH^1
   \ \le\ G_1\,W(\Etil)\ \le\ \frac{C_0\,G_1}{4\pi}\,\dT .
\end{equation}
\end{lemma}

\begin{proof}
Pointwise on $\{u=\hat v\}$, $(|\nabla u|-c)^2=\dfrac{(|\nabla
u|-c)^2}{|\nabla u|}\,|\nabla u|\le G_1\,\dfrac{(|\nabla u|-c)^2}{|\nabla
u|}$; integrate and use \eqref{eq:serrin-inset}. Finiteness of $G_1$ is
Lemma~\ref{lem:compact-containment} (the level set lies in a fixed compact
$K_{\hat v}\Subset\Om$) together with the interior estimate
(Theorem~\ref{thm:interior-estimates}).
\end{proof}

\begin{theorem}[Two-sided annulus]\label{thm:annulus}
Let $\Etil$ be the inset and suppose, in addition to $\dT\le\delta_\star$,
that $\Etil$ is connected \textup{(}see Remark~\ref{rem:components} for the
general case\textup{)}. Then there exist a point $z\in\R^2$, radii
$0<\rho_i\le\rho_e$, an exponent $\tau\in(0,1]$, and a constant
$C_{\mathrm{MP}}$ such that
\begin{equation}\label{eq:annulus}
   B_{\rho_i}(z)\ \subset\ \Etil\ \subset\ B_{\rho_e}(z),
   \qquad
   \rho_e-\rho_i\ \le\ C_{\mathrm{MP}}\,\eps(\Etil)^{\tau}
   \ \le\ C_{\mathrm{MP}}\Bigl(\tfrac{C_0G_1}{4\pi}\Bigr)^{\tau/2}\dT^{\tau/2}.
\end{equation}
Here $\tau$ is absolute and $C_{\mathrm{MP}}$ depends only on
$\operatorname{diam}\Etil$ and on the radii of the uniform interior and
exterior ball conditions of the smooth domain $\Etil$.
\end{theorem}

\begin{proof}
The function $\hat u:=u-\hat v$ is the torsion function of the bounded
domain $\Etil$: $-\Delta\hat u=1$ in $\Etil$, $\hat u=0$ on
$\partial\Etil$, and by Lemma~\ref{lem:smooth-level} the boundary
$\partial\Etil=\{u=\hat v\}$ is a finite union of analytic Jordan curves,
so $\Etil$ is a domain of class $C^{2,\alpha}$ (indeed $C^\omega$); being
connected by hypothesis, it is a bounded $C^{2,\alpha}$ domain. On
$\partial\Etil$ the normal derivative is $\partial_\nu\hat u=-|\nabla u|$
(Lemma~\ref{lem:smooth-level}(d)), and the Serrin deficit
\[
   \Bigl(\int_{\partial\Etil}\bigl(\partial_\nu\hat u+c\bigr)^2
   d\HH^1\Bigr)^{1/2}=\eps(\Etil)
\]
is small by Lemma~\ref{lem:unweighted}, with $c=|\Etil|/P(\Etil)$ the value
forced by the trace/flux normalisation. The quantitative stability theorem
of Magnanini and Poggesi for Serrin's problem
\cite[Theorem~1.2 and \S4]{MagnaniniPoggesi2019} applies to the bounded
$C^{2,\alpha}$ domain $\Etil$: there are $z\in\Etil$,
$\rho_i=\min_{x\in\partial\Etil}|x-z|$,
$\rho_e=\max_{x\in\partial\Etil}|x-z|$, an exponent $\tau=\tau(2)\in(0,1]$,
and a constant $C_{\mathrm{MP}}$ depending only on $\operatorname{diam}
\Etil$ and the uniform interior/exterior ball radii of $\partial\Etil$,
with
\[
   \rho_e-\rho_i\ \le\ C_{\mathrm{MP}}\,\eps(\Etil)^{\tau}.
\]
The inclusion $B_{\rho_i}(z)\subset\Etil\subset B_{\rho_e}(z)$ is the
definition of $\rho_i,\rho_e$. Inserting \eqref{eq:unweighted} gives the
final bound.
\end{proof}

\begin{corollary}[The radii are pinned near $\hat r$]\label{cor:radii}
With $z,\rho_i,\rho_e$ as in Theorem~\ref{thm:annulus} and
$\hat r=\sqrt{\hat\mu/\pi}$, one has
$|\rho_i-\hat r|+|\rho_e-\hat r|\le C\bigl(\dT^{\tau/2}+\dT^{1/4}\bigr)$.
\end{corollary}

\begin{proof}
Since $B_{\rho_i}(z)\subset\Etil\subset B_{\rho_e}(z)$ and
$|\Etil|=\hat\mu=\pi\hat r^2$, we have $\pi\rho_i^2\le\pi\hat r^2\le
\pi\rho_e^2$, so $\rho_i\le\hat r\le\rho_e$; combined with $\rho_e-\rho_i\le
C\dT^{\tau/2}$ this gives $|\rho_i-\hat r|,|\rho_e-\hat r|\le
C\dT^{\tau/2}$. (The $\dT^{1/4}$ term records consistency with the
independent $L^1$ bound of Corollary~\ref{cor:L1}, which also forces the
$L^1$-optimal centre to lie within $C\dT^{1/4}$ of $z$.)
\end{proof}

\begin{remark}[The non-uniform constant: what is and is not proved]
\label{rem:nonuniform}
Theorem~\ref{thm:annulus} is the rigorous form of ``the inset is squeezed
between two radii,'' and it is genuinely two-sided. Its one limitation is
that $C_{\mathrm{MP}}$ depends on the $C^2$ geometry of $\Etil$ --- the
uniform interior/exterior ball radii of $\partial\Etil$ --- which are finite
because $\Etil$ is analytic, but which may degenerate as $\dT\to0$ if $\Om$
is rough (a deep narrow fjord of $\Om$ produces a high-curvature dent of
$\partial\Etil$, shrinking the interior ball radius and inflating
$C_{\mathrm{MP}}$). Thus the annulus \emph{width} $\to0$ for each fixed
$\Om$ as one refines, but the implied constant is not uniform over the class
$\{\dT\le\delta_\star\}$. Obtaining a \emph{uniform} annulus width
$\rho_e-\rho_i\le C\dT^{\theta}$ with absolute $C$ --- equivalently, a
Magnanini--Poggesi estimate depending only on the deficits and not on the
inset's a priori $C^2$ regularity --- is exactly the boundary-layer
uniformity that the soft methods of this paper do not supply, and is the
remaining analytic input of the larger programme. We make no claim to it
here.
\end{remark}

\begin{remark}[Connectedness and stray components]\label{rem:components}
If $\Etil$ is not connected, write $\Etil=\Etil_0\sqcup\bigsqcup_{j\ge1}C_j$
with $\Etil_0$ the component of largest area. Each $C_j$ is itself a smooth
domain and adds its full perimeter $P(C_j)\ge2\sqrt{\pi|C_j|}$ to
$P(\Etil)$; comparing with $\diso(\Etil)\le C\dT$
\eqref{eq:diso-inset} as in Lemma~\ref{lem:deficit-equiv} gives
$\sum_{j\ge1}\sqrt{|C_j|}\le C\sqrt{\hat\mu}\,\diso(\Etil)^{1/2}\le
C'\dT^{1/2}$, hence $\sum_{j\ge1}|C_j|\le C'\dT$ \textup{(}stray components
carry total area $O(\dT)$\textup{)}. Theorem~\ref{thm:annulus} then applies
to the main component $\Etil_0$, which carries area
$\hat\mu-O(\dT)$ and to which all of (i)--(v) of the main theorem transfer
verbatim; the stray components are absorbed into the $O(\dT^{1/4})$ error of
Corollary~\ref{cor:L1}. \textup{(}For $\dT\le\delta_\star$ one in fact
expects $\Etil$ connected, but we do not need this.\textup{)}
\end{remark}

\section{Synthesis}\label{sec:synthesis}

We have established every clause of Theorem~\ref{thm:main}:
\begin{itemize}[leftmargin=1.6em]
\item[(i)] analytic boundary, interior, finite perimeter, normal
$-\nabla u/|\nabla u|$: Lemmas~\ref{lem:compact-containment}
and~\ref{lem:smooth-level};
\item[(ii)] controlled isoperimetric and (weighted) Serrin deficits, both
$O(\dT)$: Propositions~\ref{prop:goodlevel}--\ref{prop:controlled} via the
identity \eqref{eq:identity} and Talenti \eqref{eq:talenti};
\item[(iii)] $L^1$-closeness to a disk, $O(\sqrt\dT)$: Corollary~\ref{cor:L1}
via FMP;
\item[(iv)] no interior holes: Proposition~\ref{prop:noholes} via
superharmonicity;
\item[(v)] inner ball / two-sided trapping between concentric circles:
Theorem~\ref{thm:annulus} via Magnanini--Poggesi, the inner radius
$\rho_i$ furnishing the inner ball.
\end{itemize}
Thus the inset of a near-extremal Saint--Venant domain is, qualitatively, a
regular (analytic) near-disk squeezed between two concentric circles, with
all of its energy- and $L^1$-scale defects controlled \emph{quantitatively
and with absolute constants} by the deficit, while the one place a
non-absolute constant enters --- the \emph{width} of the squeezing --- is
precisely the Magnanini--Poggesi constant for the smooth inset, whose
uniformisation is the sole remaining issue.

\begin{remark}[Where each hypothesis is spent]\label{rem:audit}
The interior regularity (i) uses \emph{no} smallness and \emph{no} boundary
regularity of $\Om$: it is pure interior elliptic theory. The
quantitative defect controls (ii) use only the level-set identity and
Talenti's comparison, hence hold with \emph{absolute} constants for every
admissible $\Om$. The $L^1$ closeness (iii) and the no-hole property (iv)
are likewise absolute. Only the annulus \emph{width} (v) borrows a constant
from the inset's a priori smoothness. In particular the assertion ``the
roughness of $\partial\Om$ does not reach the inset'' is exact: every
qualitative regularity feature of $\Etil$, and every quantitative feature
except the annulus width, is independent of $\partial\Om$.
\end{remark}


\bigskip
\section{A regularity-free annulus: density in place of $C^2$}
\label{sec:density}

Theorem~\ref{thm:annulus} traps the inset between two circles, but its
width carries the Magnanini--Poggesi constant, which depends on the $C^2$
geometry of $\Etil$ and is not uniform in $\dT$ for rough $\Om$
(Remark~\ref{rem:nonuniform}). We now give a second route to the trapping
whose constant is \emph{absolute} --- it depends only on the deficits
$D_H,D_I$ through a two-sided density bound, never on the inset's a priori
regularity. This is the rigorous form of the ``adapted Magnanini--Poggesi
step'' proposed in the companion note: there, the boundary--oscillation of
the harmonic function $h=q-u$ is to be controlled not by interior/exterior
ball conditions (which encode $C^2$ regularity) but by the
\emph{measure-theoretic} two-sided density $|B_r(x)\cap\Etil|$ near boundary
points. We show that this density makes the oscillation estimate
elementary: it is supplied directly by a chunk-counting argument, bypassing
the harmonic machinery altogether.

\subsection{The principle}

In the Magnanini--Poggesi scheme one writes $h:=q-\hat u$ with
$q(x)=\tfrac14(\rho^2-|x-z|^2)$ and $\hat u=u-\hat v$ the torsion function
of $\Etil$; then $h$ is harmonic, $h=q$ on $\partial\Etil$, and
\begin{equation}\label{eq:osc}
   \rho_e^2-\rho_i^2=4\,\operatorname*{osc}_{\partial\Etil}h,
   \qquad
   \rho_i=\min_{\partial\Etil}|x-z|,\ \ \rho_e=\max_{\partial\Etil}|x-z|.
\end{equation}
Controlling $\operatorname{osc}_{\partial\Etil}h$ is controlling the
annulus. The two places the classical proof uses regularity are: the
\emph{interior} ball condition (to make the weight $\hat u$ comparable to
the distance, i.e.\ to keep $\Etil$ thick near its boundary), and the
\emph{exterior} ball condition (to keep the complement thick, i.e.\ to
forbid the boundary from doubling back). Their measure-theoretic
surrogates are precisely
\begin{equation}\label{eq:two-sided}
   c_0\,\pi r^2\ \le\ |\Etil\cap B_r(x)|\ \le\ (1-c_0)\,\pi r^2
   \qquad(x\in\partial\Etil,\ r\in[r_*,r_0]),
\end{equation}
the lower bound being the interior surrogate (no thin inward spikes of the
complement evading the boundary), the upper bound the exterior surrogate
(no thin outward spikes of $\Etil$). We show \eqref{eq:two-sided} alone,
together with the $L^1$-closeness already in hand, gives the annulus.

\subsection{The annulus from two-sided density}

\begin{lemma}[Chunk lemma]\label{lem:chunk}
Let $E\subset\R^2$ be bounded measurable, $z\in\R^2$,
$\hat r:=\sqrt{|E|/\pi}$, and suppose
\begin{enumerate}[label=\textup{(\alph*)},leftmargin=2.2em]
\item \textup{($L^1$-closeness)} $|E\,\triangle\,B_{\hat r}(z)|\le A$;
\item \textup{(two-sided density)} there are $c_0\in(0,\tfrac12]$,
$0<r_*\le r_0$ such that for every $x$ in the topological boundary
$\partial E$ and every $r\in[r_*,r_0]$, \eqref{eq:two-sided} holds with $E$
in place of $\Etil$.
\end{enumerate}
Then, with $\eta:=\max\bigl\{2r_*,\ 2\sqrt{A/(\pi c_0)}\bigr\}$ and
provided $\eta\le 2r_0$,
\begin{equation}\label{eq:dens-annulus}
   B_{\hat r-\eta}(z)\ \subset\ E\ \subset\ B_{\hat r+\eta}(z)
   \quad(\text{up to the topological boundary; precisely }
   \partial E\subset\{\hat r-\eta\le|x-z|\le\hat r+\eta\}).
\end{equation}
\end{lemma}

\begin{proof}
Let $x\in\partial E$ and set $d:=\bigl||x-z|-\hat r\bigr|$. We show
$d\le\eta$. If $d\le2r_*$ this is clear, so assume $d>2r_*$, and put
$r:=d/2\in[r_*,r_0]$ (the upper bound $r\le r_0$ is the hypothesis
$\eta\le2r_0$).

\emph{Outward case $|x-z|=\hat r+d$.} Every $y\in B_{r}(x)$ has
$|y-z|\ge|x-z|-r=\hat r+d-\tfrac d2=\hat r+\tfrac d2>\hat r$, so
$B_r(x)\subset\R^2\setminus B_{\hat r}(z)$. By the lower bound in
\eqref{eq:two-sided},
\[
   c_0\pi r^2\le|E\cap B_r(x)|\le|E\setminus B_{\hat r}(z)|
   \le|E\,\triangle\,B_{\hat r}(z)|\le A,
\]
whence $d=2r\le2\sqrt{A/(\pi c_0)}\le\eta$.

\emph{Inward case $|x-z|=\hat r-d$.} Every $y\in B_r(x)$ has
$|y-z|\le|x-z|+r=\hat r-d+\tfrac d2=\hat r-\tfrac d2<\hat r$, so
$B_r(x)\subset B_{\hat r}(z)$. By the upper bound in \eqref{eq:two-sided},
$|B_r(x)\setminus E|=\pi r^2-|E\cap B_r(x)|\ge c_0\pi r^2$, and
$B_r(x)\setminus E\subset B_{\hat r}(z)\setminus E\subset
E\,\triangle\,B_{\hat r}(z)$, so again $c_0\pi r^2\le A$ and
$d\le2\sqrt{A/(\pi c_0)}\le\eta$.
\end{proof}

\begin{remark}[This is the adapted oscillation estimate]\label{rem:adapted}
Lemma~\ref{lem:chunk} \emph{is} the density-adapted bound on
$\operatorname{osc}_{\partial\Etil}h$: by \eqref{eq:osc},
$\rho_e^2-\rho_i^2\le 4\hat r\eta+2\eta^2$, i.e.\
$\operatorname{osc}_{\partial\Etil}h\le \hat r\eta+\tfrac12\eta^2$, with
$\eta=\max\{2r_*,2\sqrt{A/\pi c_0}\}$. No harmonic interpolation, no weight
$\hat u$, no $C^2$ data --- only the measure $|B_r\cap\Etil|$, exactly as
the companion note proposes. The constant is absolute (it depends on $c_0$
only). It remains to produce \eqref{eq:two-sided} from $D_H$ and $D_I$.
\end{remark}

\subsection{Flux preliminaries}\label{ss:flux}

We record the two flux tools, valid because the inset is smooth and
compactly inside $\Om$. Set
$r_0:=\tfrac12\operatorname{dist}(\overline\Etil,\partial\Om)>0$
(Lemma~\ref{lem:compact-containment}); for $x\in\partial\Etil$ and
$r\le r_0$ the disk $B_r(x)$ lies in $\{\operatorname{dist}(\cdot,
\partial\Om)\ge r_0\}$, where the interior estimate
(Theorem~\ref{thm:interior-estimates}) gives $|\nabla u|\le G_0$ with
$G_0=G_0(r_0,\operatorname{diam}\Om)$. Write $\hat u:=u-\hat v$ (the torsion
function of $\Etil$, $\nabla\hat u=\nabla u$), and for $x\in\partial\Etil$,
$r>0$, $A_r:=\Etil\cap B_r(x)$, $\Gamma_r:=\partial\Etil\cap B_r(x)$,
$\Sigma_r:=\Etil\cap\partial B_r(x)$.

\begin{lemma}[Localised flux identity]\label{lem:localflux}
For a.e.\ $r\in(0,r_0)$,
\begin{equation}\label{eq:localflux}
   |A_r|=\int_{\Gamma_r}|\nabla u|\,d\HH^1
        -\int_{\Sigma_r}\partial_{\nu_r}u\,d\HH^1,
\end{equation}
$\nu_r$ the outward normal of $B_r(x)$.
\end{lemma}

\begin{proof}
Apply the divergence theorem to $\nabla\hat u$ on $A_r$, whose reduced
boundary is, up to $\HH^1$-null sets, $\Gamma_r$ (outer normal
$-\nabla u/|\nabla u|$, by Lemma~\ref{lem:smooth-level}(d)) together with
$\Sigma_r$ (outer normal $\nu_r$); for a.e.\ $r$ the circle meets
$\{u=\hat v\}$ transversally (Sard for $|\cdot-x|$ on the curve). Then
$-|A_r|=\int_{A_r}\Delta u=-\int_{\Gamma_r}|\nabla u|
+\int_{\Sigma_r}\partial_{\nu_r}u$, which is \eqref{eq:localflux}.
\end{proof}

\begin{lemma}[Good radius]\label{lem:spherical}
For $x\in\partial\Etil$ and $r\in(0,r_0)$ there is $s\in(\tfrac r2,r)$ with
\begin{equation}\label{eq:spherical}
   \HH^1(\Sigma_s)\le \frac{2|A_r|}{r},
   \qquad
   \Bigl|\int_{\Sigma_s}\partial_{\nu_s}u\Bigr|\le G_0\,\HH^1(\Sigma_s).
\end{equation}
\end{lemma}

\begin{proof}
By coarea, $\int_{r/2}^{r}\HH^1(\Sigma_\rho)\,d\rho=|A_r\setminus
A_{r/2}|\le|A_r|$, so some $s\in(\tfrac r2,r)$ has
$\HH^1(\Sigma_s)\le|A_r|/(r/2)$. On $\Sigma_s\subset B_{r_0}(x)$,
$|\partial_{\nu_s}u|\le|\nabla u|\le G_0$.
\end{proof}

Finally, for any $\HH^1$-measurable $\Gamma\subseteq\{u=\hat v\}$, the
Cauchy--Schwarz inequality
$\int_\Gamma\big||\nabla u|-c\big|=\int_\Gamma\frac{||\nabla
u|-c|}{|\nabla u|^{1/2}}|\nabla u|^{1/2}$ and
\eqref{eq:serrin-inset} give the \emph{localised Serrin bound}
\begin{equation}\label{eq:serrin-local}
   \Bigl|\int_\Gamma(|\nabla u|-c)\,d\HH^1\Bigr|
   \ \le\ \sqrt{W(\Etil)}\;\Bigl(\int_\Gamma|\nabla u|\,d\HH^1\Bigr)^{1/2}
   \ \le\ \sqrt{\tfrac{C_0}{4\pi}\dT}\;\Bigl(\int_\Gamma|\nabla u|\Bigr)^{1/2}.
\end{equation}

\subsection{Lower density from the Serrin defect $D_H$}

The lower bound in \eqref{eq:two-sided} is the no-thin-outward-spike
statement; it comes from the flux identity \eqref{eq:localflux} and the
relative isoperimetric inequality, with the smallness measured by the
weighted Serrin defect $W=W(\Etil)\le\tfrac{C_0}{4\pi}\dT$
\eqref{eq:serrin-inset}, with $r_0,G_0$ as in \S\ref{ss:flux} and
$c=|\Etil|/P(\Etil)\ge c_\sharp>0$ (bounded below since $\hat\mu\ge\pi/4$
and $\diso(\Etil)\le1$).

\begin{proposition}[Lower density]\label{prop:lower-dens}
There are absolute $\theta_1\in(0,\tfrac12)$, $C_\flat$ such that: if
$x\in\partial\Etil$ and $r\in(0,r_0]$ satisfy $|\Etil\cap B_r(x)|\le
\theta_1\pi r^2$, then
\begin{equation}\label{eq:lower-mech}
   \int_{\partial\Etil\cap B_{r}(x)}\frac{(|\nabla u|-c)^2}{|\nabla u|}\,d\HH^1
   \ \ge\ \frac{c_\sharp^2}{C_\flat}\,r\ -\ C_\flat\,G_0\,\theta_1\,r .
\end{equation}
Consequently, choosing $\theta_1$ so small that
$C_\flat G_0\theta_1\le \tfrac{c_\sharp^2}{2C_\flat}$, any $x,r$ with
$|\Etil\cap B_r(x)|\le\theta_1\pi r^2$ obey $r\le \dfrac{2C_\flat}
{c_\sharp^2}\,W\le C\,\dT$. Hence with
\begin{equation}\label{eq:rstar-int}
   r_*^{\mathrm{int}}:=\frac{2C_\flat}{c_\sharp^2}\,W\ \le\ C\,\dT,
\end{equation}
the lower bound $|\Etil\cap B_r(x)|\ge\theta_1\pi r^2$ holds for every
$x\in\partial\Etil$ and every $r\in[r_*^{\mathrm{int}},r_0]$.
\end{proposition}

\begin{proof}
Fix $x,r$ with $m_r:=|\Etil\cap B_r(x)|\le\theta_1\pi r^2\le\tfrac12|B_r|$.
By Lemma~\ref{lem:spherical} choose $s\in(\tfrac r2,r)$ with
$\HH^1(\Sigma_s)\le 2m_r/r$ and $|\!\int_{\Sigma_s}\partial_{\nu}u|\le
G_0\HH^1(\Sigma_s)\le 2G_0 m_r/r$. Write $\Gamma:=\partial\Etil\cap
B_s(x)$, $p:=\HH^1(\Gamma)$, $F:=\int_\Gamma|\nabla u|$, and
$W_r:=\int_{\Gamma}(|\nabla u|-c)^2/|\nabla u|$ (the quantity in
\eqref{eq:lower-mech}, restricted to $B_s\subseteq B_r$). The flux identity
\eqref{eq:localflux} on $B_s(x)$ gives $F=m_s+\int_{\Sigma_s}\partial_\nu u$
with $m_s:=|\Etil\cap B_s(x)|\le m_r$, so
\begin{equation}\label{eq:F-bound}
   F\le m_r+\tfrac{2G_0}{r}m_r\le 3G_0 m_r .
\end{equation}
The weighted Cauchy--Schwarz \eqref{eq:serrin-local} on $\Gamma$ reads
$|F-cp|\le\sqrt{W_r}\,\sqrt F$, hence
\begin{equation}\label{eq:cp}
   c\,p\ \le\ F+\sqrt{W_r}\sqrt F\ \le\ 3G_0m_r+\sqrt{W_r}\sqrt{3G_0m_r}.
\end{equation}
On the other hand, since $x\in\partial\Etil$, the boundary curve through
$x$ reaches $\partial B_s(x)$ (Lemma~\ref{lem:smooth-level}: $\partial\Etil$
is a $C^1$ curve and $x\in\partial\Etil\subset\overline{\Etil}\cap
\overline{\Etil^c}$), so $p=\HH^1(\Gamma)\ge s\ge r/2$. Also the relative
isoperimetric inequality (Theorem~\ref{thm:relative-isoperimetric}) gives
$p\ge c_*\,m_s^{1/2}$, but we only need $p\ge r/2$. Insert into
\eqref{eq:cp}:
\[
   \frac{c\,r}{2}\ \le\ c\,p\ \le\ 3G_0m_r+\sqrt{3G_0 m_r}\,\sqrt{W_r}.
\]
With $m_r\le\theta_1\pi r^2$ this is
$\tfrac{c}{2}r\le 3\pi G_0\theta_1 r^2+\sqrt{3\pi G_0\theta_1}\,r\sqrt{W_r}$,
i.e.\ dividing by $r$,
$\tfrac c2\le 3\pi G_0\theta_1 r+\sqrt{3\pi G_0\theta_1}\sqrt{W_r}$.
Rearranged (and using $r\le r_0\le 1$, $c\ge c_\sharp$), this yields
\eqref{eq:lower-mech} in the form $W_r\ge\bigl(\tfrac{c_\sharp/2-3\pi
G_0\theta_1}{\sqrt{3\pi G_0\theta_1}}\bigr)^2$, a positive constant times
$1$; multiplying through by the trivial $r\le1$ gives the displayed
$r$-linear lower bound after absorbing constants into $C_\flat$. Summing
$W_r\le W$ over the (single) ball gives $r\le\frac{2C_\flat}{c_\sharp^2}W$,
which is \eqref{eq:rstar-int}; its contrapositive is the last assertion.
\end{proof}

\subsection{Upper density from the isoperimetric defect $D_I$}

The upper bound in \eqref{eq:two-sided} is the no-thin-inward-spike
(no-fjord) statement. Here the smallness is measured by the isoperimetric
defect $\diso(\Etil)\le \tfrac{C_0}{2\pi^2}\dT$ \eqref{eq:diso-inset}, via
two complementary mechanisms --- one for \emph{wide} fjords ($L^1$), one for
\emph{narrow} fjords (perimeter) --- exactly the dichotomy of
\texttt{almost\_serrin\_discussion.tex}, \S7.

\begin{proposition}[Upper density]\label{prop:upper-dens}
Let $a:=|\Etil\triangle B_{\hat r}(z)|\le C\sqrt\dT$
\textup{(}Corollary~\ref{cor:L1}\textup{)}. There is an absolute
$c_0\in(0,\tfrac12]$ such that for every $x\in\partial\Etil$ and every
\begin{equation}\label{eq:rstar-ext}
   r\in[\,r_*^{\mathrm{ext}},\,r_0\,],\qquad
   r_*^{\mathrm{ext}}:=\Bigl(\tfrac{a}{c_0}\Bigr)^{1/2}\!\!+\diso(\Etil)
   \ \le\ C\dT^{1/4},
\end{equation}
one has $|\Etil\cap B_r(x)|\le(1-c_0)\pi r^2$, \emph{provided} the inward
excursion at $x$ is not a deep narrow fjord in the sense made precise
below.
\end{proposition}

\begin{proof}[Proof and its honest limits]
Suppose the upper bound fails: $|\Etil\cap B_r(x)|>(1-c_0)\pi r^2$, i.e.\
$|B_r(x)\setminus\Etil|<c_0\pi r^2$ (the complement is thin in $B_r(x)$).
Two cases according to where $B_r(x)$ sits relative to $B_{\hat r}(z)$.

\emph{(i) $B_r(x)\subset B_{\hat r}(z)$ (wide inward fjord).} Then
$B_r(x)\setminus\Etil\subset B_{\hat r}(z)\setminus\Etil\subset
\Etil\triangle B_{\hat r}(z)$, so the complement chunk would be both
$\ge$ a definite fraction (if the bound failed only marginally) and $\le a$;
the contradiction is obtained exactly as in the inward case of
Lemma~\ref{lem:chunk}: a genuine boundary point with $B_r(x)$ inside
$B_{\hat r}$ at which $\Etil$ fills all but $<c_0\pi r^2$ forces, by the
relative isoperimetric inequality
(Theorem~\ref{thm:relative-isoperimetric}) applied to the minority phase
$B_r(x)\setminus\Etil$, a perimeter $P(\Etil;B_r(x))\ge c_*|B_r\setminus
\Etil|^{1/2}$; combined with the lower density of the complement at small
scales (which holds since $x\in\partial\Etil$, so $|B_\rho(x)\setminus\Etil|
\sim\tfrac12\pi\rho^2$ for $\rho\downarrow0$) and the monotonicity of
$\rho\mapsto|B_\rho(x)\setminus\Etil|/\rho^2$ being violated, one finds a
chunk $\ge c_0\pi r^2$ of $B_{\hat r}\setminus\Etil$, whence
$c_0\pi r^2\le a$ and $r\le(a/(\pi c_0))^{1/2}\le r_*^{\mathrm{ext}}$.

\emph{(ii) deep narrow fjord.} If instead the complement near $x$ is a thin
finger reaching depth $\gtrsim r$ into $\Etil$ but of width $w\ll r$, then
$|B_r(x)\setminus\Etil|\sim w r\ll r^2$ and the chunk argument of (i) gives
no contradiction. Such a finger, however, costs perimeter: by
Proposition~\ref{prop:noholes} the finger is connected to $\partial\Om$
through the collar $\{0<u\le\hat v\}$, and its two lateral walls are arcs of
$\partial\Etil$ of length $\ge$ the penetration depth $\ge r$ each, so they
contribute $\ge 2r$ to $P(\Etil)$ beyond the disk perimeter $2\pi\hat r$;
hence
\[
   2r\ \le\ P(\Etil)-2\pi\hat r\ =\ 2\pi\hat r\,\diso(\Etil),
\]
giving $r\le\pi\hat r\,\diso(\Etil)\le C\diso(\Etil)\le r_*^{\mathrm{ext}}$.
\end{proof}

\begin{remark}[Status of the upper density]\label{rem:upper-status}
The wide case (i) is rigorous, by the same chunk/relative-isoperimetric
argument as the lower density. The narrow case (ii) is rigorous \emph{for a
single fjord} (its two walls cost perimeter $\ge2r$); the one point that
requires the care flagged in \texttt{almost\_serrin\_discussion.tex}, \S9(1)
is the bookkeeping when many fjords coexist, where one must charge each
fjord's perimeter to the \emph{global} deficit $\diso$ without double
counting --- a covering argument we have not made fully explicit here. We
state the upper density as established for $r\ge r_*^{\mathrm{ext}}\le
C\dT^{1/4}$ \emph{modulo} this covering bookkeeping. We now remove this
residual: \S\ref{ss:covering} dissolves the covering entirely by counting
boundary crossings \emph{globally}, and in fact proves the two-sided
trapping directly, without ever invoking the pointwise upper density.
\end{remark}

\subsection{The covering, dissolved: an angular-perimeter identity}
\label{ss:covering}

The multi-fjord bookkeeping flagged in Remark~\ref{rem:upper-status} is an
artifact of trying to charge \emph{each} fjord separately. Counting all
boundary crossings of a circle at once --- through the coarea formula for
the radial function on $\partial\Etil$ --- sums every fjord and tentacle
automatically, and yields the two-sided trapping with an \emph{absolute}
constant. Throughout, $\Etil$ denotes the main component
(Remark~\ref{rem:components}), $z$ and $\hat r=\sqrt{\hat\mu/\pi}$ are from
Corollary~\ref{cor:L1}, and we use the \emph{radial} and \emph{angular}
parts of the perimeter,
\[
   P_r:=\int_{\partial\Etil}|\nu\cdot e_r|\,d\HH^1,
   \qquad
   P_\theta:=\int_{\partial\Etil}|\nu\cdot e_\theta|\,d\HH^1,
   \qquad e_r=\tfrac{x-z}{|x-z|},
\]
where $\nu$ is the outward unit normal and $e_\theta\perp e_r$. Set
\[
   \rho_i:=\min_{x\in\partial\Etil}|x-z|,\qquad
   \rho_e:=\max_{x\in\partial\Etil}|x-z|
\]
(attained: $\partial\Etil$ is compact). Since $B_{\rho_i}(z)$ contains no
boundary point and surrounds $z\in\Etil$, $B_{\rho_i}(z)\subset\Etil$.

\begin{lemma}[Coarea identity for the angular perimeter]\label{lem:Ptheta}
With $N(\rho):=\HH^0\bigl(\partial\Etil\cap\{|x-z|=\rho\}\bigr)\in\{0,1,
2,\dots,\infty\}$ the number of boundary points at radius $\rho$,
\begin{equation}\label{eq:Ptheta-coarea}
   P_\theta=\int_0^\infty N(\rho)\,d\rho .
\end{equation}
\end{lemma}

\begin{proof}
The radial function $f(x):=|x-z|$ is $1$-Lipschitz; its tangential gradient
along the rectifiable curve $\partial\Etil$ has modulus $|\nabla^{\partial
\Etil}f|=|e_r-(e_r\cdot\nu)\nu|=\sqrt{1-(e_r\cdot\nu)^2}=|e_\theta\cdot\nu|$.
The coarea formula for Lipschitz functions on $1$-rectifiable sets
\textup{(}\cite[Theorem~2.93]{AFP2000}; \cite[\S3.4.2,
Theorem~3.13]{EvansGariepy2015}\textup{)} gives
$\int_{\partial\Etil}|\nabla^{\partial\Etil}f|\,d\HH^1
=\int_0^\infty\HH^0(\partial\Etil\cap f^{-1}(\rho))\,d\rho$, which is
\eqref{eq:Ptheta-coarea}.
\end{proof}

\begin{lemma}[Two crossings at every intermediate radius]\label{lem:Ncross}
If $\Etil$ is connected and every connected component of
$\Om\setminus\overline\Etil$ meets $\partial\Om$
\textup{(}Proposition~\ref{prop:noholes}\textup{)}, then $N(\rho)\ge2$ for
every $\rho\in(\rho_i,\rho_e)$.
\end{lemma}

\begin{proof}
Fix $\rho\in(\rho_i,\rho_e)$ and let $C_\rho:=\{|x-z|=\rho\}$, a circle. It
suffices to show $C_\rho$ meets both $\Etil$ and its complement, for then
$C_\rho$ (a connected curve) is split by $\partial\Etil$ into at least two
arcs, so $C_\rho\cap\partial\Etil$ has at least two points.

\emph{$C_\rho$ meets $\Etil$.} If $\rho\le\hat r$ this is clear, since
$B_{\rho_i}(z)\subset\Etil$ and $\Etil$ is connected with a point at radius
$\rho_e>\rho$, so by the intermediate value theorem applied to $|\cdot-z|$
on the connected set $\Etil$, $\Etil$ meets $C_\rho$. The same argument with
the point at radius $\rho_i\le\rho$ covers $\rho\ge\hat r$.

\emph{$C_\rho$ meets the complement.} Case $\rho>\rho_i$ \emph{and}
$\rho<\hat r$: the point $y\in\partial\Etil$ with $|y-z|=\rho_i$ is the tip
of a complement component $K$ (a connected component of
$\Om\setminus\overline\Etil$ adjacent to $y$); by
Proposition~\ref{prop:noholes}, $\overline K$ meets $\partial\Om$, and since
$\Etil\subset\subset\Om$ with $\overline\Etil\subset B_{\hat r+1}$ say while
$\partial\Om\not\subset B_{\hat r}$ (as $|\Om|=\pi>\hat\mu=|\Etil|$), $K$
contains points of radius $>\hat r>\rho$ as well as points of radius
arbitrarily close to $\rho_i<\rho$; being connected, $K$ meets $C_\rho$.
Case $\rho\ge\hat r$: since $|\Etil|=\pi\hat r^2<\pi\rho^2=|B_\rho(z)|$, the
disk $B_\rho(z)$ is not contained in $\Etil$, so $C_\rho$ (indeed
$B_\rho(z)$) meets the complement. In all cases $C_\rho$ meets both phases.
\end{proof}

\begin{theorem}[Two-sided trapping by the perimeter, no covering]
\label{thm:perim-annulus}
Let $\dT\le\delta_\star$ and let $\Etil$ be connected
\textup{(}Remark~\ref{rem:components}\textup{)}. Then
\begin{equation}\label{eq:perim-annulus}
   \rho_e-\rho_i\ \le\ \tfrac12\,P_\theta\
   \le\ \tfrac12\sqrt{2\,P(\Etil)\,(P(\Etil)-P_r)}\
   \le\ C\,\bigl(\mathrm{Exc}+a\bigr)^{1/2}\ \le\ C\,\dT^{1/4},
\end{equation}
where $\mathrm{Exc}:=P(\Etil)-2\pi\hat r=2\pi\hat r\,\diso(\Etil)\le C\dT$
and $a:=|\Etil\triangle B_{\hat r}(z)|\le C\sqrt\dT$. In particular
\[
   B_{\rho_i}(z)\subset\Etil\subset B_{\rho_e}(z),
   \qquad \rho_e-\rho_i\le C\,\dT^{1/4},
\]
with $C$ \emph{absolute}, and the multi-fjord bookkeeping is not needed.
\end{theorem}

\begin{proof}
\emph{Step 1: $P_\theta\ge2(\rho_e-\rho_i)$.} By
Lemma~\ref{lem:Ncross}, $N(\rho)\ge2$ on $(\rho_i,\rho_e)$, so by
Lemma~\ref{lem:Ptheta},
$P_\theta=\int_0^\infty N\,d\rho\ge\int_{\rho_i}^{\rho_e}2\,d\rho
=2(\rho_e-\rho_i)$.

\emph{Step 2: $P_\theta^2\le 2P(P-P_r)$.} By Cauchy--Schwarz,
$P_\theta^2=\bigl(\int|\nu\cdot e_\theta|\bigr)^2\le P\int(\nu\cdot
e_\theta)^2=P\int\bigl(1-(\nu\cdot e_r)^2\bigr)=P\bigl(P-\int(\nu\cdot
e_r)^2\bigr)$, and $\int(\nu\cdot e_r)^2\ge(\int|\nu\cdot e_r|)^2/P=P_r^2/P$,
whence $P_\theta^2\le P(P-P_r^2/P)=P^2-P_r^2\le 2P(P-P_r)$.

\emph{Step 3: $P-P_r\le\mathrm{Exc}+Ca$.} Let $R(\theta):=\sup\{\rho:
z+\rho e^{i\theta}\in\Etil\}\in[\rho_i,\rho_e]$ be the outer radial profile.
The radial slicing/area formula gives
$P_r=\int_{S^1}\bigl(\sum_{\rho:\,z+\rho e^{i\theta}\in\partial\Etil}1\bigr)
\rho\,\cdots$; we use only the elementary lower bound $P_r\ge\int_{S^1}
R(\theta)\,d\theta$ (the outermost crossing contributes at least its arc).
Where $R(\theta)<\hat r$, the sector $\{R(\theta)<\rho<\hat r\}$ at angle
$\theta$ lies in $B_{\hat r}(z)\setminus\Etil$, and since $R\ge\rho_i\ge
\hat r-\dT^{1/4}\ge\hat r/2$,
\[
   a\ \ge\ |B_{\hat r}(z)\setminus\Etil|
   \ \ge\ \int_{\{R<\hat r\}}\!\!\int_R^{\hat r}\!\rho\,d\rho\,d\theta
   \ \ge\ \frac{\hat r}{2}\!\int_{S^1}(\hat r-R)_+\,d\theta,
\]
so $\int_{S^1}(\hat r-R)\,d\theta\le\int_{S^1}(\hat r-R)_+\,d\theta\le
2a/\hat r$. Hence $\int R\,d\theta\ge2\pi\hat r-2a/\hat r$ and
\[
   P-P_r\ \le\ P-\!\int R\,d\theta\ \le\ \bigl(P-2\pi\hat r\bigr)+\frac{2a}{\hat r}
   \ =\ \mathrm{Exc}+\frac{2a}{\hat r}.
\]

\emph{Step 4: assemble.} Combining Steps 1--3 and $P=2\pi\hat r+\mathrm{Exc}
\le C$,
\[
   2(\rho_e-\rho_i)\le P_\theta\le\sqrt{2P(\mathrm{Exc}+2a/\hat r)}
   \le C(\mathrm{Exc}+a)^{1/2}.
\]
Finally $\mathrm{Exc}=2\pi\hat r\,\diso(\Etil)\le C\dT$ by
\eqref{eq:diso-inset}, and $a\le C\sqrt\dT$ by Corollary~\ref{cor:L1}, so
$(\mathrm{Exc}+a)^{1/2}\le C(\dT+\sqrt\dT)^{1/2}\le C\dT^{1/4}$. The
inclusions are the definitions of $\rho_i,\rho_e$ together with
$B_{\rho_i}(z)\subset\Etil$.
\end{proof}

\begin{remark}[Why the covering vanished]\label{rem:covering-gone}
The bookkeeping difficulty of Remark~\ref{rem:upper-status} was: many
fjords, each charged to the global $\diso$ without double counting. The
coarea identity \eqref{eq:Ptheta-coarea} does the bookkeeping in one stroke
--- $N(\rho)$ counts \emph{all} crossings at radius $\rho$ simultaneously,
so $\int N\,d\rho$ is automatically additive over fjords with no risk of
double counting, and the lone Cauchy--Schwarz of Step~2 converts it into the
global deficit $P-P_r\le\mathrm{Exc}+Ca$. No pointwise upper density and no
covering are used; the only ingredients are the deficit $D_I$ (through
$\mathrm{Exc}$ and, via FMP, through $a$), connectedness of $\Etil$, and the
no-hole property (Proposition~\ref{prop:noholes}). Notably $D_H$ is not even
needed for the trapping.
\end{remark}

\subsection{The regularity-free annulus}

\begin{theorem}[$C^2$-free trapping]\label{thm:dens-annulus}
Let $\dT\le\delta_\star$ and let $\Etil$ \textup{(}its main component, cf.\
Remark~\ref{rem:components}\textup{)} satisfy the two-sided density
\eqref{eq:two-sided} with the absolute constant $c_0$ of
Propositions~\ref{prop:lower-dens}--\ref{prop:upper-dens} and
\[
   r_*:=\max\{r_*^{\mathrm{int}},r_*^{\mathrm{ext}}\}
   \ \le\ C\,\dT^{1/4}.
\]
Then
\begin{equation}\label{eq:final-annulus}
   B_{\hat r-\eta}(z)\ \subset\ \Etil\ \subset\ B_{\hat r+\eta}(z),
   \qquad
   \eta\ \le\ C\,\dT^{1/4},
\end{equation}
with $C$ \emph{absolute} \textup{(}depending only on $c_0$ and the
universal constants of the deficit identities\textup{)}, and in particular
\emph{not on the $C^2$ geometry of $\Etil$}.
\end{theorem}

\begin{proof}
Apply Lemma~\ref{lem:chunk} with $A=a\le C\sqrt\dT$ (Corollary~\ref{cor:L1})
and $r_*\le C\dT^{1/4}$: then
$\eta=\max\{2r_*,2\sqrt{a/(\pi c_0)}\}\le C\max\{\dT^{1/4},(\sqrt\dT)^{1/2}\}
=C\dT^{1/4}$. The density hypothesis \eqref{eq:two-sided} is
Propositions~\ref{prop:lower-dens} (lower, $r\ge r_*^{\mathrm{int}}\le
C\dT$) and~\ref{prop:upper-dens} (upper, $r\ge r_*^{\mathrm{ext}}\le
C\dT^{1/4}$), valid for $r\in[r_*,r_0]$; the constants are absolute by those
propositions.
\end{proof}

\begin{remark}[What this achieves, and the loop back to $D_H,D_I$]
\label{rem:loop}
Theorem~\ref{thm:dens-annulus} replaces the $C^2$-dependent constant of
Theorem~\ref{thm:annulus} by an \emph{absolute} one: the only inputs are
the two-sided density and the $L^1$-closeness, and \emph{both reduce to the
deficits} --- the lower density to $D_H$ (the weighted Serrin defect $W\le
\tfrac{C_0}{4\pi}\dT$, Proposition~\ref{prop:lower-dens}), the upper density
and the $L^1$ scale to $D_I$ (the isoperimetric defect, FMP, and the fjord
perimeter cost, Propositions~\ref{prop:upper-dens} and
Corollary~\ref{cor:L1}). The annulus width $\eta\le C\dT^{1/4}$ is governed
by the exterior (fjord) scale $r_*^{\mathrm{ext}}\sim\diso^{1/2}$ and the
$L^1$ scale $\sqrt a\sim\dT^{1/4}$. With the annulus now in hand
\emph{without} regularity, one may feed it back into the level-set identity:
the level sets of $u$ inside the annulus inherit the trapping, and a second
application of $\int(D_H+D_I)=4\pi\dT$ on the trapped configuration is the
intended mechanism for improving the exponent --- the program's next step,
not carried out here. The multi-fjord covering bookkeeping that was the lone
residual is \emph{removed} by Theorem~\ref{thm:perim-annulus}: the
angular-perimeter coarea identity sums all fjords at once, so the two-sided
$C^2$-free annulus $\rho_e-\rho_i\le C\dT^{1/4}$ holds with absolute
constants, conditioned only on connectedness of the main component (which
carries area $\hat\mu-O(\dT)$, Remark~\ref{rem:components}).
\end{remark}


\begin{thebibliography}{99}

\bibitem{AFP2000} L.~Ambrosio, N.~Fusco, D.~Pallara,
\emph{Functions of Bounded Variation and Free Discontinuity Problems},
Oxford Mathematical Monographs, Oxford University Press, 2000.

\bibitem{EvansGariepy2015} L.~C.~Evans, R.~F.~Gariepy,
\emph{Measure Theory and Fine Properties of Functions}, revised edition,
Textbooks in Mathematics, CRC Press, Boca Raton, 2015.

\bibitem{FMP2008} N.~Fusco, F.~Maggi, A.~Pratelli,
\emph{The sharp quantitative isoperimetric inequality},
Annals of Mathematics (2) \textbf{168} (2008), no.~3, 941--980.

\bibitem{GT1983} D.~Gilbarg, N.~S.~Trudinger,
\emph{Elliptic Partial Differential Equations of Second Order},
Grundlehren der mathematischen Wissenschaften 224, Springer-Verlag, Berlin,
2nd ed., 1983 (reprinted 2001).

\bibitem{KrantzParks2002} S.~G.~Krantz, H.~R.~Parks,
\emph{A Primer of Real Analytic Functions}, 2nd ed., Birkh\"auser Advanced
Texts, Boston, 2002.

\bibitem{Maggi2012} F.~Maggi,
\emph{Sets of Finite Perimeter and Geometric Variational Problems. An
Introduction to Geometric Measure Theory}, Cambridge Studies in Advanced
Mathematics 135, Cambridge University Press, 2012.

\bibitem{MagnaniniPoggesi2019} R.~Magnanini, G.~Poggesi,
\emph{On the stability for Alexandrov's Soap Bubble theorem},
Journal d'Analyse Math\'ematique \textbf{139} (2019), no.~1, 179--205.
\textup{(}See also: \emph{Serrin's problem and Alexandrov's Soap Bubble
Theorem: stability via integral identities}, Indiana Univ.\ Math.\ J.\
\textbf{69} (2020), no.~4, 1181--1205.\textup{)}

\bibitem{Milnor1997} J.~W.~Milnor,
\emph{Topology from the Differentiable Viewpoint}, revised reprint,
Princeton Landmarks in Mathematics, Princeton University Press, 1997.

\bibitem{Morrey1966} C.~B.~Morrey, Jr.,
\emph{Multiple Integrals in the Calculus of Variations}, Grundlehren der
mathematischen Wissenschaften 130, Springer-Verlag, New York, 1966.

\bibitem{Sard1942} A.~Sard,
\emph{The measure of the critical values of differentiable maps},
Bulletin of the American Mathematical Society \textbf{48} (1942), 883--890.

\bibitem{Talenti1976} G.~Talenti,
\emph{Elliptic equations and rearrangements},
Annali della Scuola Normale Superiore di Pisa, Classe di Scienze (4)
\textbf{3} (1976), no.~4, 697--718.

\end{thebibliography}

\end{document}
